% ============================================================================
% PLOS ONE submission draft
% Title: How cells decide: thermodynamics and hierarchical signal control
%        in Bacillus subtilis sporulation
% Author: Eugênio Simão
% Status: working draft (June 2026)
%
% This file is a self-contained article-class manuscript that mirrors the
% PLOS ONE structure (Abstract -> Author Summary -> Introduction -> Methods
% -> Results -> Discussion -> References).  When ready for submission, the
% body can be transferred into the official PLOS LaTeX template
% (plos_latex_template.tex / plos2015.bst).
% ============================================================================

\documentclass[11pt,a4paper]{article}

\usepackage[utf8]{inputenc}
\usepackage[T1]{fontenc}
\usepackage[english]{babel}
\usepackage{amsmath,amssymb,amsthm}
\usepackage{booktabs}
\usepackage{geometry}
\usepackage{graphicx}
\usepackage{hyperref}
\usepackage{cleveref}
\usepackage{microtype}
\usepackage{xcolor}
\usepackage[font=small,labelfont=bf]{caption}
\usepackage{setspace}   % double spacing (PLOS ONE)
\usepackage[right]{lineno} % line numbers (PLOS ONE review)

\geometry{margin=2.5cm}
\hypersetup{colorlinks=true,linkcolor=blue!50!black,citecolor=blue!50!black,urlcolor=blue!50!black}
% PLOS ONE: abbreviate Figure to Fig in cross-references
\crefname{figure}{Fig}{Figs}
\Crefname{figure}{Fig}{Figs}

\newcommand{\kB}{k_{\mathrm{B}}}
\newcommand{\Gs}{G_s}
\newcommand{\Ge}{G_E}
\newcommand{\thetaeff}{\theta_{\mathrm{eff}}}
\newcommand{\DeltaF}{\Delta F}
\newcommand{\Aset}{\mathcal{A}}
\newcommand{\BC}{\mathrm{BC}}
\newcommand{\CV}{\mathrm{CV}}
\newcommand{\Hfate}{H_{\mathrm{fate}}}

% ─────────────────────────────────────────────────────────────────────────────
\begin{document}
\doublespacing   % PLOS ONE requires double spacing
\linenumbers     % PLOS ONE requires continuous line numbers
\vspace*{0.2in}

\begin{flushleft}
{\Large
\textbf{Non-Equilibrium Thermodynamics of Biological Decision-Making:
A Formal Framework and \textit{Bacillus subtilis} Validation}
}
\newline
\\
Eug\^{e}nio Sim\~{a}o$^{1,*}$
\\
\bigskip
\textbf{1} Department of Computer Science, Universidade Federal de Santa
Catarina (UFSC), Ararangu\'a, Santa Catarina, 88906-072, Brazil
\\
\bigskip
* eugenio.simao@ufsc.br
\end{flushleft}

\begin{abstract}
\noindent
How does a single cell commit irreversibly to a fate?
In this work we argue that cell-fate commitment is an irreversible
thermodynamic transition in an open system --- a controlled expense of
free energy that reduces the cell's uncertainty about its own future ---
and we introduce Signal Hierarchical Petri Nets (SHPN), a formal
modelling language whose thermodynamic content is readable directly
from network topology rather than inferred from fitted rate constants.
We validate the framework on \textit{Bacillus subtilis} sporulation,
driven by three open questions from Fujita and Losick (2005):
(Q1)~why does abrupt Spo0A* induction yield $\approx 5\%$ sporulation
while the gradual phosphorelay achieves $\approx 52\%$?
(Q2)~why is $\sigma_H$ necessary even when Spo0A$\sim$P is abundant?
(Q3)~what determines the $\approx 50\%$ bet-hedging fraction in isogenic cells?
Five numerical experiments and three targeted biological sweeps answer
these thermodynamically.
(Q1)~The key variable is the nutrient medium: SM medium depletes ATP
below the vegetative-basin floor ($2.13$~mM), triggering commitment
regardless of the route; CH medium keeps ATP above the floor, rendering
both routes ineffective.
(Q2)~Zeroing $\sigma_H$ transcription abolishes sporulation entirely
($0\%$) even with an abrupt Spo0A$\sim$P pulse; $\sigma_H$ concentrates
$>96\%$ of circuit entropy production and is the thermodynamic
commitment gate, not a downstream reporter.
(Q3)~Non-mean-field critical exponents ($\gamma_{\mathrm{sub}} \approx 2.4$,
$\gamma_{\mathrm{sup}} \approx 4.6$) place the bifurcation in a noise-driven
universality class, making $\approx 50\%$ the thermodynamically natural
outcome without parameter steering.
Additional findings: cumulative dissipated work (${\sim}10^{8}\,\kB T$)
renders equilibrium fluctuation theorems inapplicable --- the
signature of genuine biological irreversibility.
The framework is applicable to any commitment circuit for which the
regulatory topology is known.
\end{abstract}

\bigskip

% ─────────────────────────────────────────────────────────────────────────────
\section*{Author Summary}
\addcontentsline{toc}{section}{Author Summary}

\noindent
Cells commit to a fate --- spore, neuron, immune effector --- by spending energy
to reduce uncertainty about their own future.  Once committed, the path back costs
more than any living cell can afford.  How much does that decision cost?  What
makes it sharp or gradual, permanent or reversible?

We present Signal Hierarchical Petri Nets (SHPN), a modelling language that
encodes the causal ordering and energy cost of each commitment step directly in
the network topology.  We apply it to \textit{Bacillus subtilis} sporulation
and address three classical experimental observations: why does an abrupt pulse
of the master regulator Spo0A fail where a gradual phosphorelay succeeds;
why is the sigma factor $\sigma_H$ a necessary gate for commitment;
and why do exactly $\approx 50\%$ of genetically identical cells sporulate
in nutrient-limited medium.

Three lessons emerge.  First, \emph{deciding is cheap}: the signalling
cascade that commits the cell uses less than 4\% of the circuit's total
energy budget.  Second, \emph{executing is expensive}: building the spore
coat dominates the energy bill by more than an order of magnitude.
Third, the key variable in the Spo0A experiment is the nutrient medium:
depleted medium drops ATP below the commitment threshold; rich medium does not.

\bigskip
\hrule
\bigskip

% ─────────────────────────────────────────────────────────────────────────────
\section{Introduction}
\label{sec:intro}

\subsection{A Decision Without a Decider}

A bacterium runs out of food.  Within an hour it will either keep
running the long-shot lottery of vegetative growth --- hoping nutrients
return before its envelope fails --- or commit, irreversibly, to
building a spore: a near-crystalline, metabolically silent capsule that
can wait centuries for better weather.  The two paths are not
reconcilable.  A cell midway through coat assembly cannot retract the
decision; a cell that has bet on growth cannot recover the half-spore
it never started.

And yet there is no decider.  No homunculus weighs the options.  The
cell is a few cubic micrometres of crowded chemistry, and the choice
emerges from the interaction of a handful of phosphorylated kinases
with their substrates, modulated by the noise inherent in counting
molecules whose copy numbers fall in the dozens.  The decision is real
--- it is measurable, repeatable, and consequential --- yet it is made
by no one in particular.  This is the central puzzle of cell-fate
commitment, and a version of it appears wherever life faces a fork:
in the stem cell choosing a tissue lineage, in the lymphocyte choosing
between clonal expansion and anergy, in the algal cell choosing between
photosynthesis and predation.

Waddington, intuiting the problem long before the mechanisms were
known, pictured development as a ball rolling on a tilted landscape of
valleys and ridges~\cite{Waddington1957}.  The metaphor is good --- so
good that it has survived seventy years of molecular biology --- but
it is, on close inspection, a picture of an effect rather than an
explanation of a cause.  What does the landscape mean physically?
Where does the tilt come from?  Why do valleys exist at all?  And, the
question this paper centres on: \emph{what does the cell actually
spend, and on what, to roll one way rather than the other?}

\subsection{Three Questions from Fujita and Losick (2005)}
\label{ssec:fujita-questions}

In 2005, Fujita and Losick published a decisive set of experiments on
\textit{B.~subtilis} sporulation that sharpened the puzzle to its
essential form~\cite{Fujita2005}.  Their findings reduce to three
explicit questions that a physical theory of commitment must answer ---
and that any adequate model must reproduce \emph{without parameter steering}.

A critical but unremarked feature of the Fujita design is the \emph{medium
asymmetry}: the two routes were tested in completely different carbon
environments.  The natural phosphorelay experiment was carried out in
\textbf{SM medium} (Sterlini-Mandelstam sporulation medium, no carbon
source except glutamate) --- nutrient-depleted conditions under which
the cell's ATP pool declines as catabolism exhausts available substrate.
The abrupt Spo0A* experiment was carried out in
\textbf{CH medium} (casein hydrolysate, a rich carbon source) ---
the cell was actively growing in exponential phase (OD$_{600} = 0.4$)
when Spo0A* synthesis was induced by IPTG.
Fujita did not measure ATP in either experiment; the medium difference
was reported as a technical necessity of the inducible-promoter
strategy, not as a variable under investigation.
Our SHPN model predicts that this medium asymmetry is \emph{mechanistically
essential}: the vegetative-basin ATP floor ($2.13$~mM, Sweep~C)
is the thermodynamic boundary that SM medium allows the cell to cross
while CH medium does not.  The rich medium experiment is not merely
a genetic perturbation; it is an energetic trap.

\textbf{Q1 (timing and medium asymmetry).}
Abrupt induction of constitutively active Spo0A* in \textbf{CH medium}
(rich, exponential-phase cells) yielded only $\approx 5\%$ sporulation
(Fig.~2C/D), while gradual KinA induction in \textbf{SM medium}
(nutrient-limited cells) achieved $\approx 52\%$ (Fig.~2A/B).
\emph{Why does the combination of route and medium matter?
Is it the abruptness of Spo0A activation, the richness of the medium,
or both?}

\textbf{Q2 (gate necessity).}
$\sigma_H$-null mutants failed to sporulate even when Spo0A* was
induced abruptly (Fig.~4; $0\%$ forespore formation).
\emph{Why is $\sigma_H$ a necessary condition for commitment if it is
downstream of the apparent decision-maker Spo0A$\sim$P?}

\textbf{Q3 (bet-hedging fraction).}
Under the natural gradual route only $\approx 52\%$ of isogenic cells
sporulated in the same SM medium --- not $100\%$, not $0\%$.
\emph{What mechanism produces this specific, reproducible partial
commitment across a genetically uniform population?}

Classical biological formalisms describe these observations but cannot
explain them from first principles: the $52\%/5\%$ ratio and the
$\sigma_H$ requirement are inserted as parameters, not derived from
circuit topology.  A physical theory should derive --- or at minimum
bound --- all three from the thermodynamic structure of the circuit,
with no outcome steering.  The five numerical experiments in
\Cref{sec:results} each address one or more of these questions; the
explicit answers are compiled in \Cref{tab:covalidation}.

\subsection{Physics as the Missing Foundation}

Living cells are not at thermodynamic equilibrium.  They cannot be:
equilibrium is the definition of death.  A living cell continuously
imports free energy --- principally as the phosphoanhydride bonds of
ATP and GTP --- and exports entropy, in the form of heat and chemical
waste, faster than it generates it internally.  This puts cells
squarely in the domain of \emph{non-equilibrium thermodynamics of open
systems}~\cite{Nicolis1977}, a discipline whose central
claim is that ordered, organised behaviour is possible in nature only
because some structures are willing to pay for it in dissipated free
energy.

From this perspective a fate commitment is not just a transition
between cell states.  It is the cell purchasing certainty about its
own future, at a price set by the laws of physics: by Landauer's
principle~\cite{Landauer1961}, by Schnakenberg's accounting of
chemical-network dissipation~\cite{Schnakenberg1976}, by Wang's
construction of an effective landscape from the stationary distribution
of an irreversible Markov chain~\cite{Wang2011}, and by the celebrated
fluctuation theorems of Jarzynski~\cite{Jarzynski1997} and
Crooks~\cite{Crooks1999} that bound the heat dissipated by any process
relative to its free-energy difference.  Each of these tools, taken
alone, says something concrete about cells.  Taken together they begin
to describe what kind of physical event a fate commitment actually is.

This is a strong base on which to stand, but it is a base without a
building: the physics gives us the right vocabulary, yet none of the
modelling formalisms used by biologists today \emph{speaks} that
vocabulary natively.  Ordinary differential equations express
concentrations; Boolean networks express logic gates; stochastic Petri
nets express discrete firings.  None of them expresses, as part of its
syntax, the thermodynamic content of a commitment.  To compute the
physical quantities the framework demands, one must first build the
model, then post-process trajectories, then hope that one's bookkeeping
is consistent with one's theoretical claims.  Something is missing.

Important work has applied the non-equilibrium physics vocabulary
\emph{to} biological models post hoc: Ge and Qian derived exact
relations between dissipation and landscape curvature for chemical
networks~\cite{GeQian2010}; Ao constructed potential landscapes for
stochastic differential equations~\cite{Ao2004}; Rao and Esposito gave
a complete stochastic thermodynamics for chemical reaction
networks~\cite{RaoEsposito2016}.  What none of these contributions
provides, however, is a \emph{modelling language} --- a syntax in which
the thermodynamic structure of a commitment is encoded at the time the
model is built, rather than extracted by analysis after the fact.  The
distinction matters: when thermodynamics is a post-processing step,
nothing in the model prevents topological choices that are
thermodynamically incoherent, and nothing in the formalism flags the
inconsistency.

\subsection{What This Paper Adds}

The missing piece, we argue, is a formal modelling language whose
\emph{architecture} reflects the architecture of biological
decision-making, and whose syntax has intrinsic thermodynamic meaning.
We call this language Signal Hierarchical Petri Nets (SHPN), and develop it in
\Cref{sec:formalism}.  It extends standard stochastic Petri nets in
three directions, each of which corresponds to a property of
biological commitment systems that has been observed empirically for
decades but never given a syntactic home: the separation of
\emph{execution} from \emph{control}; the strict \emph{causal
ordering} of signal resolution; and the parameterisation of each
commitment step by its own thermodynamic profile.

Three concrete contributions follow.  First, in \Cref{sec:physics} we
establish the non-equilibrium thermodynamic vocabulary needed to
characterise biological commitment as a physical event, and identify
five measurable quantities that any adequate commitment model must
compute.  Second, in
\Cref{sec:formalism} we introduce the formal modelling language SHPN,
whose architecture is designed so that each of those five quantities is
computable directly from the topology of the model.  Third, in
\Cref{sec:results}, we validate the framework against the
\textit{B.~subtilis} sporulation cascade through five complementary
numerical experiments and three targeted biological sweeps that answer
the Fujita questions directly.

The result, we hope, is the beginning of a unified physical picture:
cells decide by paying free energy to collapse uncertainty about their
own future, and the architecture of that decision can be made explicit
enough to be modelled, simulated, and tested.  \Cref{sec:physics}
establishes the thermodynamic framework; \Cref{sec:formalism} presents
SHPN; \Cref{sec:methods,sec:results} describe the \textit{B.~subtilis}
validation; \Cref{sec:discussion} interprets the findings and situates
the framework among alternatives; \Cref{sec:conclusion} closes.

% ─────────────────────────────────────────────────────────────────────────────
\section{Thermodynamic Framework for Biological Commitment}
\label{sec:physics}

\subsection{Cells Are Open Systems}

A living cell maintains steep concentration gradients, pumps ions against
electrochemical potentials, and consumes ATP at rates orders of magnitude
above equilibrium.  Equilibrium thermodynamics is therefore inapplicable;
the appropriate framework is non-equilibrium thermodynamics of open
systems~\cite{Nicolis1977}.  The second law holds globally,
but a cell locally exports entropy --- as heat and chemical waste --- faster
than it generates it, sustaining itself as a \emph{dissipative structure}.
A cell that achieves equilibrium with its environment is dead, by
definition.

\subsection{Commitment Reduces Uncertainty at a Thermodynamic Cost}

A pluripotent cell carries uncertainty about its future fate: many
phenotypes are accessible.  Let \(\Aset(\Ge)\) denote the set of
quasi-stationary attractors of the cell's biochemical network --- the
phenotypic possibility space encoded in the genotype.  Under a uniform
prior, the initial fate entropy is
\begin{equation}
  \Hfate^{(0)} = \ln|\Aset(\Ge)| \quad\text{nats}.
  \label{eq:fate-entropy}
\end{equation}
For a binary commitment system such as sporulation
(\(|\Aset(\Ge)| = 2\)), \(\Hfate^{(0)} = \ln 2 \approx 0.69\) nats
(equivalently 1 bit).  Commitment collapses
this distribution onto a single attractor.  By Landauer's
principle~\cite{Landauer1961} the minimum free energy cost of this
erasure is
\begin{equation}
  \DeltaF_{\mathrm{commitment}} \;\geq\; \kB T \cdot \Hfate^{(0)}
  \;=\; \kB T \ln 2 \;\approx\; 2.97\times10^{-21}\,\mathrm{J}
  \quad(T = 310.15\,\mathrm{K}).
  \label{eq:landauer-binary}
\end{equation}
This is a floor, not a prediction.  How close any real or simulated
commitment system approaches the bound is an empirical question; the
ratio
\begin{equation}
  \eta_L = \frac{\kB T \cdot \Hfate^{(0)}}{\DeltaF_{\mathrm{measured}}}
  \;\leq\; 1
  \label{eq:eta-L}
\end{equation}
is the \emph{Landauer efficiency} of the commitment process.

\subsection{Entropy Production Rate Decomposes by Reaction}

For an open chemical reaction network the Schnakenberg entropy production
rate~\cite{Schnakenberg1976} is
\begin{equation}
  \dot\sigma \;=\;
  \sum_\rho (J^+_\rho - J^-_\rho)\,\ln\!\frac{J^+_\rho}{J^-_\rho},
  \label{eq:schnakenberg}
\end{equation}
the sum of contributions from each reaction \(\rho\) with forward and
reverse fluxes \(J^\pm_\rho\).  For unidirectional driven reactions
(NTP-hydrolysis-coupled transitions) the contribution is well-approximated
by \(\bar J_\rho \cdot \Delta G_\rho\), where \(\bar J_\rho\) is the mean
firing rate and \(\Delta G_\rho\) the free energy released per firing.
The Schnakenberg formula is one of the few exact statements available for
chemical networks far from equilibrium; the decomposition into
per-transition contributions is the foundation of NET~5.
Note that pure phosphotransfer reactions (no NTP hydrolysis) have
$J^+_\rho \approx J^-_\rho$ near biochemical equilibrium, so their
individual contributions to $\dot\sigma$ are negligible; in NET~5
they are reported as $\dot\sigma_\rho = 0$ to reflect this approximation.

\subsection{The Non-Equilibrium Landscape}

Wang and collaborators~\cite{Wang2011} extended the Waddington picture
beyond equilibrium by defining the effective landscape as
\begin{equation}
  U_{\mathrm{eff}}(M) = -\ln P_{\mathrm{ss}}(M),
  \label{eq:noneq-landscape}
\end{equation}
where \(P_{\mathrm{ss}}\) is the stationary probability of marking \(M\)
under the full stochastic dynamics.  The decomposition of the
probability current into a gradient (landscape) and rotational (flux) term
is exact, but computing \(P_{\mathrm{ss}}\) for a high-dimensional gene
circuit requires solving the stationary Fokker--Planck equation in full
dimension --- intractable for any realistic model.  An empirical surrogate
--- the kernel-density estimate of marker markings across replicates ---
is feasible and forms the basis of NET~2 (\Cref{fig:waddington-basin}).

\begin{figure}[!ht]
  \centering
  \includegraphics[width=\linewidth]{figures/fig_waddington_shypn.png}
  \caption{\textbf{Epigenetic potential landscape of \textit{B.~subtilis}
  sporulation.}
  $U_{\mathrm{eff}}(\varphi, t) = -\ln P_{\mathrm{ss}}(\varphi, t)$
  computed from 20 stochastic replicates of the SHPN model (Methods).
  The order parameter $\varphi$ integrates phosphorelay, $\sigma$-factor,
  and structural-place trajectories (see NET~2, Results).
  White line: mean $\varphi(t)$.
  Yellow circle: commitment crossing at $t = 16.8$\,min
  ($\theta_{\mathrm{eff}}$ threshold).
  Green star: first mature spore at $t = 20.2$\,min.
  Red dashed line: $\theta_{\mathrm{eff}}$ time.}
  \label{fig:waddington-basin}
\end{figure}

\subsection{Fluctuation Theorems Break for Strongly Driven Cascades}

The Jarzynski equality~\cite{Jarzynski1997} and Crooks fluctuation
theorem~\cite{Crooks1999} require Markovian dynamics, a single thermal
bath, and dissipated work \(W\) on the order of \(\kB T\).  For
multi-stage commitment cascades, independent estimates of mean ATP
consumption per sporulation cycle ($\sim 10^7$ molecules, each
releasing $\sim 23\,\kB T$) predict \(W \sim 2\times10^{8}\,\kB T\) before any
simulation is run.  This is a \emph{prediction from the known biology},
not a design choice of the model's causal ordering.  NET~3 confirms it:
the exponential average underlying the Jarzynski equality underflows in
finite-precision arithmetic for every accessible trajectory --- the
numerical signature of the predicted dissipation scale.  The failure is
not a bug introduced by the POSet architecture but a direct measurement
of how far from equilibrium the cascade operates; detailed balance is not
even approximately recoverable.

\subsection{Critical Behaviour at the Commitment Bifurcation}

Near a commitment bifurcation parameterised by a control variable
\(N\) (e.g. ambient nutrient level), the variance of the order parameter
diverges as
\begin{equation}
  \CV^2 \sim C\,|N - N_c|^{-\gamma},
  \label{eq:critical-scaling}
\end{equation}
where the exponent \(\gamma\) characterises the universality class.  A
mean-field saddle-node gives \(\gamma = 1\); departures from \(\gamma=1\)
indicate non-mean-field behaviour driven by stochastic fluctuations
rather than gradient flow.  NET~4 measures \(\gamma\) on each side of
\(N_c\).

\subsection{Summary: Five Measurable NESM Properties of a Commitment Circuit}

A commitment circuit is fully characterised, thermodynamically, by:
(i)~its Landauer efficiency \(\eta_L\) (\Cref{eq:eta-L});
(ii)~the number and depth of attractors in \(U_{\mathrm{eff}}\)
(\Cref{eq:noneq-landscape});
(iii)~the dissipated work magnitude \(W/\kB T\) relative to fluctuation
theorem applicability;
(iv)~the critical exponents \(\gamma_{\pm}\) on each side of the
bifurcation (\Cref{eq:critical-scaling});
(v)~the Schnakenberg EPR decomposition by reaction
(\Cref{eq:schnakenberg}).
What is required for these to be computable from a model is a formalism
whose topology directly encodes them.  This is the role of SHPN.

% ─────────────────────────────────────────────────────────────────────────────
\section{The Formalism: Signal Hierarchical Petri Nets (SHPN)}
\label{sec:formalism}

\subsection{Standard SPN Background}

A standard stochastic Petri net is a bipartite directed graph
\((P, T, F)\) with places \(P\) (state variables, holding discrete
\emph{tokens}), transitions \(T\) (events that consume and produce
tokens), and arcs \(F \subseteq (P\times T) \cup (T\times P)\) labelled
by integer weights.  The marking \(M : P \to \mathbb{N}\) is the
state.  A transition \(t\) fires when each input place \(p\) has at
least the weight of the arc \((p,t)\) tokens; firing consumes the
input weights and produces the output weights.  In the
\emph{stochastic} variant, firing follows Poisson statistics with a
rate function \(\Phi_t(M)\)~\cite{Marsan1995,Wilkinson2018}.

\subsection{SHPN as a 13-Tuple}

SHPN extends the standard SPN tuple by adding signal places, signal
flow arcs with Hill-like parameters, a layer function, and an
enablement operator that checks the layer constraint:
\begin{equation}
  \mathrm{SHPN} \;=\; (P, T, F, W, M_0, \Phi, C, F_t, \Psi, F_s, W_s,
  \lambda, \theta).
  \label{eq:13-tuple}
\end{equation}
The classical components are \(P,T,F,W,M_0,\Phi,C,F_t\): places,
transitions, normal arcs, weights, initial marking, rate functions,
capacities, and test arcs (\(F_t\): presence-checking arcs that do not
consume tokens).  The novel components are:
\begin{itemize}
\item \(\Psi \subseteq P\): \emph{signal places} (regulatory, non-mass-bearing).
\item \(F_s \subseteq (\Psi \times T) \cup (T \times \Psi)\): \emph{signal
  flow arcs}, each carrying a Hill-like parameter tuple
  \(\Gamma = (K,n,\varepsilon)\); these arcs consume tokens with weight
  \(W_s : F_s \to \mathbb{R}^{+}\), the signal arc weight function.
\item \(\lambda : \Psi \to \mathbb{N}_0\): a layer function inducing a
  POSet \(\preceq\) on the signal places.  The layer of a transition
  is \(\ell(t) = \max\{\lambda(p) : (p,t) \in F_s\}\) (the highest
  layer among its signal inputs; 0 if none).
\item \(\theta\): per-transition effective thresholds derived from
  \(\Gamma\) (see Primitive~3).
\end{itemize}

\subsection{Three Architectural Primitives}

The three architectural primitives of SHPN are: (i)~the factorisation
of the cell's state space into a \emph{genotype manifold} and a
\emph{signal manifold}; (ii)~a partial order on the signal manifold,
enforced as a syntactic constraint; and (iii)~a thermodynamic gate
parameter tuple on each signal arc.  The first of these is the
conceptual heart of the formalism and deserves its own treatment.

% ─── Primitive 1 ─────────────────────────────────────────────────────────────
\subsubsection{Primitive~1 --- Two State Spaces: Genotype and Signal}
\label{sssec:two-spaces}

Every SHPN model partitions its places into two disjoint subsets,
\(P = P_E \sqcup \Psi\), where \(P_E\) carries mass-bearing biochemical
species (metabolites, proteins, mRNAs, structural assemblies) and
\(\Psi\) carries regulatory signals (phosphorylation states, active
\(\sigma\) factors, environmental cues).  The graph factorises as
\begin{equation}
  \mathrm{SHPN} \;=\; \Ge \,\times\, \Gs
  \quad\text{with}\quad
  \Ge = (P_E, T, F, \Phi), \qquad
  \Gs = (\Psi, F_s, \lambda, \Gamma),
  \label{eq:factorisation}
\end{equation}
and the full marking space splits into a product
\(\mathcal{M} = \mathcal{M}_E \times \mathcal{M}_s\),
where \(\mathcal{M}_E = \mathbb{N}^{|P_E|}\) is the \emph{genotype state
space} and \(\mathcal{M}_s = \mathbb{N}^{|\Psi|}\) is the \emph{signal state space}.

The execution subgraph \(\Ge\) is fixed by the cell's genome:
stoichiometry, network topology, and rate functions \(\Phi\) do not
change during a commitment event.  What \(\Ge\) determines is the set
of quasi-stationary attractors
\begin{equation}
  \Aset(\Ge) \;=\; \bigl\{\, m_E^{*} \in \mathcal{M}_E :
  m_E^{*} \text{ is a quasi-stationary state of } \Ge \,\bigr\},
  \label{eq:attractor-set}
\end{equation}
the phenotypic possibility space.
For \textit{B.~subtilis} sporulation \(|\Aset(\Ge)| = 2\) (vegetative
cell and mature spore); for a haematopoietic progenitor it may be eight
or ten.  No phenotype is created during commitment --- every reachable
fate is already implicit in \(\Aset(\Ge)\) before the decision begins.
The role of the signal cascade is exclusively to \emph{select} among
pre-existing possibilities.

The signal subgraph \(\Gs\) does not add attractors.  It progressively
biases the probability distribution over \(\Aset(\Ge)\), draining mass
from some attractors and concentrating it on others.  The effective
non-equilibrium landscape (\Cref{eq:noneq-landscape}) becomes
conditioned on the signal marking:
\begin{equation}
  U_{\mathrm{eff}}(m_E \mid m_s) \;=\; -\ln P_{\mathrm{ss}}(m_E \mid m_s),
  \label{eq:tilted-landscape}
\end{equation}
where \(m_s\) is the \emph{tilt} that reshapes the landscape as the
cascade progresses.  The resulting genotype-to-phenotype map is
\begin{equation}
  \mathrm{phenotype} \;=\;
  \arg\min_{\,m_E^{*} \,\in\, \Aset(\Ge)} \;
  U_{\mathrm{eff}}\!\bigl(m_E^{*} \,\bigm|\, m_s^{\mathrm{final}}\bigr):
  \label{eq:argmin}
\end{equation}
genotype provides the space of possibilities; signal traces a path through it.

Phenotype selection is entropy collapse, and it has a price.  Under a
uniform prior over \(\Aset(\Ge)\) the cell's fate entropy is
\begin{equation}
  \Hfate^{(0)} \;=\; \ln\bigl|\Aset(\Ge)\bigr|.
  \label{eq:H0}
\end{equation}
Commitment drives this to zero by concentrating all probability mass
on a single attractor.  By Landauer's principle the minimum free-energy
cost of this collapse is
\begin{equation}
  \boxed{\;\DeltaF_{\mathrm{development}} \;\geq\;
  \kB T \ln\bigl|\Aset(\Ge)\bigr|.\;}
  \label{eq:dev-cost}
\end{equation}
An organism with a richer fate repertoire pays proportionally more;
the cost of multipotency is built into the thermodynamic accounting.

\Cref{tab:two-spaces} schematises the two-state-space decomposition.
Between commitment events the two spaces evolve nearly independently;
at a commitment gate they couple through the \(\Gamma\) parameters of
the signal flow arc, mutual information between \(m_E\) and \(m_s\)
becomes nonzero, and the Landauer cost is paid in dissipated ATP/GTP
within \(\Ge\).  A notable consequence: because the signal cascade
is stochastic (\(\Gs\) is driven by Poisson-sampled transitions),
isogenic cells will trace different paths through \(\mathcal{M}_s\) and
land on different attractors of \(\Aset(\Ge)\).  Bet-hedging is therefore
not an evolved strategy bolted on top of a deterministic module, but a
thermodynamic inevitability written into the factorisation
\eqref{eq:factorisation} whenever \(|\Aset(\Ge)| > 1\).

\begin{table}[h]
\centering
\caption{The two state spaces composing every SHPN model and their
respective biological and thermodynamic semantics.  The \(\Gamma\)
tuples on signal flow arcs are the bridge that couples them at
commitment events.}
\label{tab:two-spaces}
{\small
\begin{tabular}{@{}lp{4.2cm}p{4.2cm}@{}}
\toprule
                       & \textbf{Genotype space} \(\mathcal{M}_E\)
                       & \textbf{Signal space} \(\mathcal{M}_s\) \\
\midrule
Subgraph               & \(\Ge = (P_E, T, F, \Phi)\)
                       & \(\Gs = (\Psi, F_s, \lambda, \Gamma)\) \\
Elements               & Metabolites, proteins, structures
                       & Active kinases, \(\sigma\) factors, second messengers \\
Encodes                & What the cell \emph{can become}
                       & What the cell \emph{is becoming} \\
Persistence            & Genomically fixed; reusable across contexts
                       & Run-specific; reset between experiments \\
Currency               & Free energy (J, mol, \(\Delta G\))
                       & Information (bits, decisions resolved) \\
Quasi-stationary       & \(\Aset(\Ge)\): the phenotype repertoire
                       & Transient trajectories through layers \(\lambda\) \\
Initial entropy        & \(\Hfate^{(0)} = \ln|\Aset(\Ge)|\) nats
                       & Maximally uncommitted at \(t=0\) \\
End-state entropy      & \(\approx 0\) after commitment
                       & Locked at \(m_s^{\mathrm{final}}\) \\
\bottomrule
\end{tabular}}
\end{table}

\medskip\noindent\textit{Why this matters: an epistemic asymmetry.}
The genotype is shared by every cell of the same species ---
information-theoretically, it is \emph{public}.  The phenotype, by
contrast, is the outcome of an individual cell's specific signal
history: which environmental cues it encountered, in which order,
with which stochastic fluctuations.  Two genetically identical cells in
the same medium can land on different attractors of \(\Aset(\Ge)\) ---
not because their genomes differ, but because their trajectories
through \(\mathcal{M}_s\) did.  The conditional entropy
\begin{equation}
  H(\mathrm{phenotype} \mid \mathrm{genotype}) \;=\; \Hfate^{(0)} \;>\; 0
  \qquad \text{whenever } |\Aset(\Ge)| > 1
\end{equation}
is therefore not a defect of the model --- it is the formal expression
of population-level heterogeneity.  In any system with \(|\Aset(\Ge)|
> 1\) and stochastic dynamics on \(\Gs\), isogenic cells \emph{must}
diversify across fates; bet-hedging is not an evolved strategy bolted
on top of a deterministic decision module, but a thermodynamic
inevitability written into the factorisation \eqref{eq:factorisation}.

% ─── Primitive 2 ─────────────────────────────────────────────────────────────
\subsubsection{Primitive~2 --- POSet Causal Ordering Enforced by PreemptionCheck}
A transition \(t\) at layer \(\ell(t)\) is enabled iff
\begin{align}
  \mathrm{Enabled}(t,M) \;\iff\;\,
  &\mathrm{NormalEnabled}(t,M) \;\wedge\;
  \mathrm{TestEnabled}(t,M) \;\wedge\; \nonumber\\
  &\mathrm{SignalEnabled}(t,M) \;\wedge\;
  \mathrm{PreemptionCheck}(t,M),
  \label{eq:enabled}
\end{align}
where PreemptionCheck verifies that every signal place at strictly lower
layer has reached its threshold:
\begin{equation}
  \mathrm{PC}(t,M) \;\iff\;
  \bigwedge_{\substack{p \in \Psi \\ \lambda(p) < \ell(t)}}
  \bigl[ M(p) \geq \thetaeff(p \to t) \bigr].
  \label{eq:preemption}
\end{equation}
This is the lattice meet over all lower-layer constraints
(\(\mathrm{PC} \equiv \mathrm{PreemptionCheck}\)), making
causal ordering a structural rather than dynamic property: a topology
that violates causality cannot fire, regardless of rate constants.

% ─── Primitive 3 ──────────────────────────────────────────────────────────────────
\subsubsection{Primitive~3 --- The \(\Gamma = (K, n, \varepsilon)\) Gate Tuple}
Each signal flow arc carries three parameters:
\begin{itemize}
  \item \(K\) --- the commitment threshold (Hill midpoint / EC\(_{50}\)):
    the concentration at which the gate fires with probability \(1/2\).
    This is precisely the maximum-entropy point of the binary commitment
    variable and locates the separatrix on the landscape.
  \item \(n\) --- the cooperativity (Hill coefficient): how sharply the
    gate switches.  Width of the high-entropy decision window scales as
    \(K/\sqrt{n}\); large \(n\) compresses the Landauer step.
  \item \(\varepsilon\) --- the reversibility floor (Marcus reorganisation
    energy~\cite{Marcus1993}): the minimum residual signal after firing;
    bounds the reverse-crossing probability through the Boltzmann factor
    \(e^{-(\theta - \varepsilon)/\kB T}\).
\end{itemize}
The effective threshold is derived analytically from \(\Gamma\):
\begin{equation}
  \thetaeff = K\Bigl(\frac{\varepsilon}{1-\varepsilon}\Bigr)^{1/n}.
  \label{eq:theta-eff}
\end{equation}
\(\Gamma\) is the complete specification of one Landauer step in
biologically measurable units, and is what makes the architectural
primitives empirically grounded rather than abstract.

\subsection{Hybrid Stochastic + ODE Dynamics}

SHPN supports four transition types: \texttt{stochastic} (\(\tau\)-leaping
Poisson sampling, used for low-copy decision steps), \texttt{continuous}
(ODE-evaluated, used for high-throughput metabolism and gene expression),
\texttt{immediate} (zero-delay rule), and \texttt{adaptive} (regime-switching).
Operator splitting per time step \(dt\) is: continuous flow update
\(\to\) stochastic \(\tau\)-leap \(\to\) finalise.  The GPU hybrid engine
(\Cref{sec:methods}) batches both legs across replicates.

% ─────────────────────────────────────────────────────────────────────────────
\section{Methods}
\label{sec:methods}

\subsection{\textit{B.~subtilis} Sporulation Model}

The validation model
\mbox{\texttt{bacillus\_sporulation\_v9.shy}}
encodes the phosphorelay-driven commitment cascade of \textit{Bacillus
subtilis} sporulation in the SHPN formalism: 41~places
(31 signal places $\Psi$ + 10 protocol parameter places $\square$,
including \texttt{k\_sigmaH\_factor} for the $\sigma_H$-gate sweep),
36~transitions, and 114~arcs.  The four signal hierarchy layers are
the starvation sensor KinA$_{\mathrm{kinase}}$ ($\lambda = 0$),
the phosphorelay integrator Spo0A$\sim$P ($\lambda = 1$),
the commitment gate $\sigma_H$ ($\lambda = 2$, $\theta_{\sigma H} = 1.60\,\mu$M),
and the execution programme $\sigma_F/\sigma_E/\sigma_G/\sigma_K$ ($\lambda = 3$).
The \texttt{k\_sigmaH\_factor} parameter place (M$_0 = 1$) scales the
$\sigma_H$ transcription rate and enables the true $\sigma_H$-null sweep
(Fujita Protocol sweep~F3).  All biological parameter values derive from
published sources; none were adjusted to match sporulation fractions.
\textbf{Note:} the five NET thermodynamic tests (\Cref{tab:sweeps}) use the
predecessor v7 model (35~places, 28~transitions) and are being re-run
against the v9 model; those results will be updated upon completion.
Network topology follows the experimentally characterised
KinA/Spo0F/Spo0B/Spo0A phosphorelay~\cite{Burbulys1991},
the SinR/SinI gate~\cite{Fujita2005},
and the hierarchical \(\sigma\)-factor programme
(\(\sigma_H, \sigma_F, \sigma_E, \sigma_G, \sigma_K\)) documented by
Stragier and Losick~\cite{Stragier1996} and Errington~\cite{Errington2003};
initial phosphorelay species concentrations follow fluorescence
measurements by Fujita and Losick~\cite{Fujita2005}; rate constants
for ATP-regenerating and sigma-factor transcription transitions are
derived from the BRENDA enzyme kinetics database~\cite{BRENDA}.
The token unit is \(\mu\)M; physical constants throughout are
\(T = 310.15\,\mathrm{K}\), cell volume \(V = 1\,\mathrm{fL}\),
\(\Delta G_{\mathrm{ATP}} = \Delta G_{\mathrm{GTP}} = 60\,\mathrm{kJ/mol}
= 23.27\,\kB T\) per molecule.

\begin{sloppypar}
The decision layer (six \texttt{stochastic} transitions) covers
\texttt{T\_KinA\_activation}, \texttt{T\_Spo0F\_phosphorylation},
\texttt{T\_Spo0F\_dephos}, \texttt{T\_Spo0A\_\allowbreak phosphorylation},
\texttt{T\_Spo0A\_dephosphorylation}, and \texttt{T\_septation}.  The
execution layer (12 \texttt{continuous} transitions) covers
\(\sigma\)-factor transcription cascades (\(\sigma_H\), \(\sigma_F\),
\(\sigma_E\), \(\sigma_G\), \(\sigma_K\)) and the structural assembly
programme (forespore formation, cortex, inner/outer coat, spore
maturation).  Remaining transitions handle metabolic bookkeeping
(ATP/GTP regeneration, nutrient depletion, SinR/SinI dynamics).
\end{sloppypar}

\begin{figure}[!ht]
  \centering
  \includegraphics[width=\linewidth]{figures/bacillus_sporulation_v7_thesis-1.png}
  \caption{\textbf{SHPN model of \textit{B.~subtilis} sporulation.}
  Circles: biological places (species concentrations, $\mu$M).
  Rectangles: transitions (all shown in black).
  Hexagons: signal places ($\Psi$).
  Arc styles: solid black = normal (consuming); dashed blue = test
  (non-consuming presence check); dashed light grey = signal-flow
  (consuming, cascade-gated); solid black with filled circle head =
  inhibitor (non-consuming, disables on presence).
  The decision layer ($G_s$, upper left) comprises six \texttt{stochastic}
  transitions encoding the KinA/Spo0F/Spo0A phosphorelay and septation.
  The execution layer ($G_E$, remainder) comprises 12 \texttt{continuous}
  transitions for $\sigma$-factor transcription cascades and structural
  assembly.}
  \label{fig:model-topology}
\end{figure}

\subsection{Simulation Engine}

Simulations used the SHyPN simulation platform's hybrid stochastic/ODE engine with GPU
acceleration.  At each macroscopic step \(dt\), the engine batches
\(N_{\mathrm{rep}}\) replicates in parallel: ODE integration of
\texttt{continuous} transitions is followed by \(\tau\)-leaping Poisson
sampling of \texttt{stochastic} transitions
(\(\tau_{\max} = 0.01\,\mathrm{s}\) throughout).  Per-transition
firing counts are accumulated without per-step trajectory storage,
enabling the Schnakenberg decomposition of NET~5.

\subsection{Sweep Designs}

Two families of sweeps were performed.  \Cref{tab:sweeps} describes the
five thermodynamic NET tests; \Cref{tab:fujita_sweeps} describes the four
Fujita Protocol sweeps addressing Q1--Q3 directly.

\begin{table}[ht]
\centering
\caption{\textbf{NET thermodynamic tests} (v9 model, 41~places;
  run identifiers archived in repository).
  All sweeps: $\tau_{\max} = 0.1$~s,
  6-h horizon, $T = 310.15$~K.  $N_{\mathrm{rep}}$: replicates per condition.}
\label{tab:sweeps}
{\small
\begin{tabular}{@{}llcc@{}}
\toprule
Test & Swept parameter & Cond. & $N_{\mathrm{rep}}$ \\
\midrule
NET~1 & $\sigma$-factor half-life (\texttt{SIGMA\_HALFLIFE\_MIN}) & 17 & 100 \\
NET~2 & Initial nutrients (\texttt{INITIAL\_NUTRIENTS}) & 13 & 100 \\
NET~3 & Nutrients step $N_0\!\to\!N_1$ at $t = 3600$~s ($2\times2$) & 4 & 200 \\
NET~4 & Initial nutrients, fine grid & 14 & 300 \\
NET~5 & Initial nutrients, Schnakenberg conditions & 3 & 300 \\
\bottomrule
\end{tabular}}
\end{table}

\begin{table}[ht]
\centering
\caption{\textbf{Fujita Protocol sweeps} (v9 model, 41~places).
  All sweeps use \texttt{k\_sigmaH\_factor}~$= 1$ and
  \texttt{LOADING\_DOSE}~$= 0$ unless stated.  $N_{\mathrm{rep}}$:
  replicates per condition.}
\label{tab:fujita_sweeps}
{\small
\begin{tabular}{@{}llp{6.0cm}cc@{}}
\toprule
Sweep & Question & Swept parameters & Cond. & $N_{\mathrm{rep}}$ \\
\midrule
F1 & Q1 & \texttt{INITIAL\_NUTRIENTS} $\times$ \texttt{LOADING\_DOSE}
    ($\{100,1440\}\times\{0,27\}$~\textmu M, $2\times2$) & 4 & 100 \\
F2 & Q1 (ext.) & \texttt{LOADING\_DOSE}~$= 27$~\textmu M fixed;
    \texttt{INITIAL\_NUTRIENTS}: 100--2200~\textmu M & 10 & 100 \\
F3 & Q2 & \texttt{k\_sigmaH\_factor} $\times$ \texttt{LOADING\_DOSE}
    ($\{0,0.2,\ldots,1\}\times\{0,27\}$, $6\times2$) & 12 & 50 \\
F4 & Q3 & \texttt{INITIAL\_NUTRIENTS}: 100--3000~\textmu M,
    natural route & 16 & 200 \\
\bottomrule
\end{tabular}}
\end{table}

\subsection{Analysis Pipelines}

Each test has a dedicated analysis script archived in the public
repository (see Reproducibility below).
The five pipelines implement: trajectory ATP-delta integration
(NET~1); Gaussian KDE of final markings and bimodality coefficient
(NET~2); forward/reverse protocol-work distributions (NET~3);
power-law fitting of \(\CV^2\) on each side of \(N_c\) (NET~4);
and per-transition Schnakenberg EPR via
\(\bar J_\rho \Delta G_\rho\) (NET~5).

\subsection{Reproducibility}

All materials --- model file, sweep configurations, analysis scripts,
and per-replicate output --- are archived under version control in the
public SHyPN repository
(\url{https://github.com/simao-eugenio/shypn},
DOI:~\href{https://doi.org/10.5281/zenodo.20835073}{10.5281/zenodo.20835073}).
Each archived run contains a model snapshot (the exact bytes loaded
by the simulation worker), a provenance record (repository state,
model checksum, timestamps), and the full per-replicate output,
allowing independent reproduction of every figure and table.

% ─────────────────────────────────────────────────────────────────────────────
\section{Results}
\label{sec:results}

\subsection{NET~1 --- Landauer Efficiency}
\label{ssec:t1}

For each replicate, gross ATP consumed is the sum of negative ATP-pool
trajectory deltas; observed dissipation is converted to joules using
\(N_A\), \(V\), and \(\Delta G_{\mathrm{ATP}}\), and the Landauer
efficiency \eqref{eq:eta-L} is evaluated against
\(\DeltaF_{\mathrm{Landauer}} = \kB T \ln 2\).
Results are shown in \Cref{tab:t1}.

\begin{table}[h]
\centering
\caption{NET~1 results: Landauer efficiency \(\eta_L\) at selected
\(\sigma\)-factor half-life conditions (100 replicates per condition;
\(V = 1\)~fL).  \(N_{\mathrm{ATP}}\): mean ATP consumed per
committing replicate.  Excess~\((\times) = 1/\eta_L\): fold-excess of observed
dissipation over the Landauer bound \(\kB T \ln 2\).
Only conditions with $>0\%$ sporulation at $N_0 = 1440\,\mu$M are shown.
Run: \texttt{run\_20260704\_181236} (v9, 100~reps/condition).}
\label{tab:t1}
\begin{tabular}{rrrrrr}
\toprule
HL (min) & Spor.\,\% & \(N_{\mathrm{ATP}}\) (\(\mu\)M) & \(\eta_L\)
  & Excess (\(\times\)) \\
\midrule
10  &   8 & 2770 & \(1.88\times10^{-8}\) & \(5.32\times10^{7}\) \\
20  & 100 & 2776 & \(1.88\times10^{-8}\) & \(5.33\times10^{7}\) \\
30  & 100 & 2794 & \(1.86\times10^{-8}\) & \(5.37\times10^{7}\) \\
60  & 100 & 2858 & \(1.82\times10^{-8}\) & \(5.49\times10^{7}\) \\
90  & 100 & 2852 & \(1.83\times10^{-8}\) & \(5.48\times10^{7}\) \\
120 &  38 & 2860 & \(1.82\times10^{-8}\) & \(5.50\times10^{7}\) \\
\bottomrule
\end{tabular}
\end{table}

\(\eta_L\) is $\approx 1.82$--$1.88 \times 10^{-8}$ across all conditions
where sporulation occurs --- biological switches are nowhere near the Landauer
bound, consistent with the large gap between the bound and the cost of real
cellular computation and with its direct experimental
verification~\cite{MehtaSchwab2012,Berut2012}.
The key finding is that $\eta_L$ is \emph{topology-constant}: ATP consumption
per committing cell varies by less than $3\%$ over the entire 10--120~min
half-life range ($2770$--$2860\,\mu$M), confirming that the energy cost
is fixed by the circuit topology, not the kinetic parameters.

However, the sporulation fraction is sharply non-monotonic in HL:
\begin{itemize}
\item HL~$= 10$~min: only $8\%$ sporulation --- $\sigma_H$ decays before
  sustaining the commitment pulse.
\item HL~$= 20$--$90$~min: $100\%$ sporulation --- optimal transient
  amplifier window.
\item HL~$= 120$~min (wild-type): $38\%$ at $N_0 = 1440\,\mu$M ---
  consistent with the F4 result ($39\%$) and confirming self-consistency.
\item HL~$\geq 150$~min: $0\%$ sporulation --- $\sigma_H$ accumulates
  to $1.3$--$5.6\,\mu$M (far above $\theta_{\sigma H} = 1.60\,\mu$M)
  but SinR never drops below $7.5\,\mu$M because the SinI synthesis
  pathway is disrupted by the altered $\sigma_H$ dynamics.
\end{itemize}
This non-monotonic window identifies the wild-type HL~$= 120$~min as
thermodynamically optimal \emph{for bet-hedging}, not for reliability.
Shorter half-lives (20--90~min) produce more reliable commitment at
$N_0 = 1440\,\mu$M but destroy the stochastic diversity required for
population-level bet-hedging.  The $\sigma_H$ half-life is therefore
\emph{not} tuned to maximise $\eta_L$ (which is topology-constant) but
to position the circuit at the maximum-entropy operating point of the
noise-driven bifurcation landscape, exactly at $N_c \approx 1370\,\mu$M.

\subsubsection*{Trajectory analysis: two mechanistically distinct commitment routes}

G5-tier per-replicate trajectories reveal two qualitatively different
commitment modes depending on half-life (\Cref{tab:t1_traj}).

\begin{table}[h]
\centering
\caption{Commitment timing from G5 trajectories (NET~1, sporulating
replicates only).  $t_{\mathrm{SinR}}$: first time SinR $< 7.5\,\mu$M.
$t_{\mathrm{sept}}$: first Septum$\,> 0$.
$[\sigma_H]_{\max}$: peak $\sigma_H$ across trajectory.}
\label{tab:t1_traj}
{\small\begin{tabular}{rcrrcr}
\toprule
HL (min) & Spor.\,\% & $t_{\mathrm{SinR}}$ (min) & $t_{\mathrm{sept}}$ (min)
  & Lag (min) & $[\sigma_H]_{\max}$ ($\mu$M) \\
\midrule
20  & 100 & $0.1 \pm 0.0$   & $246 \pm 4$   & $246$ & 0.265 \\
30  & 100 & $0.1 \pm 0.0$   & $245 \pm 3$   & $245$ & 0.393 \\
60  & 100 & $0.2 \pm 0.0$   & $245 \pm 3$   & $245$ & 0.761 \\
90  & 100 & $0.6 \pm 0.0$   & $244 \pm 3$   & $244$ & 1.113 \\
120 &  38 & $335 \pm 14$    & $335 \pm 14$  & $0.3$ & 1.476 \\
\bottomrule
\end{tabular}}
\end{table}

Two commitment modes emerge:

\textbf{Mode~A --- passive SinR drift (HL $\leq 90$~min).}
SinR falls below $7.5\,\mu$M within $0.1$--$0.6$~min via natural
continuous degradation, without SinI-mediated sequestration.
Septum firing is then delayed $\approx 244$~min until $\sigma_H$
accumulates enough for T\textsubscript{septation} to fire stochastically
at low propensity ($[\sigma_H] \approx 0.27$--$1.11\,\mu$M).  The
commitment clock is set by a nutrient-depletion timer upstream of
$\sigma_H$: despite $[\sigma_H]_{\max}$ varying $4\times$ across
HL~$= 20$--$90$~min, T\textsubscript{septation} fires at the same time
($244$--$246 \pm 3$~min) in all conditions.  The $\sigma_H$ concentration
modulates the \emph{probability} of firing during each SinR excursion
below $7.5\,\mu$M, not the \emph{timing} of the excursion.

\textbf{Mode~B --- active SinI/SinR (HL $= 120$~min, wild-type).}
SinR does not drop below $7.5\,\mu$M until $t = 335$~min, when
stochastic SinI synthesis from Spo0A\textsubscript{$\sim$P}
sequesters SinR below the inhibitor threshold.  Septum then fires
within $0.3 \pm 0.3$~min --- essentially instantaneously.
This is the classical stochastic SinI/SinR bet-hedging mechanism.
Only $38\%$ of replicates achieve the SinI burst (the Poisson-distributed
SinI synthesis event), producing the bet-hedging fraction directly.

A prediction emerges from this separation: \emph{deleting SinI} should
abolish Mode~B but leave Mode~A intact, shifting the sporulation
fraction at $N_0 = 1440\,\mu$M from $38\%$ to a value determined
solely by the passive SinR drift dynamics.  This prediction is testable
by setting SinI~$= 0$ in the sweep and comparing the septum timing
distribution.

\textbf{Paradox at HL $\geq 150$~min: silent dead-end.}
$\sigma_H$ accumulates above $\theta_{\sigma H} = 1.60\,\mu$M
(peak $1.80$--$3.35\,\mu$M) --- T\textsubscript{septation} is fully
enabled on the $\sigma_H$ axis --- yet $0\%$ sporulation.
SinR never drops below $7.5\,\mu$M, evidently because Spo0A$\sim$P
accumulation is suppressed under high-HL dynamics, blocking SinI
synthesis.  The cell enters a \emph{$\sigma_H$-locked vegetative state}:
the downstream execution programme is primed but the SinI/SinR trigger
never fires.  This constitutes a new architectural prediction --- $\sigma_H$
over-stabilisation creates a commitment-incompetent state despite
satisfying the transcriptional threshold criterion.

\subsection{NET~2 --- Non-Equilibrium Landscape Reconstruction}
\label{ssec:t2}

The empirical landscape \(U_{\mathrm{eff}}(x) = -\ln \hat P(x)\) is estimated
from the distribution of final ATP-pool values across 100 replicates per
condition (\texttt{run\_20260704\_201120}, v9 model, 13 conditions,
$N_0 = 100$--$3000\,\mu$M, LOADING\_DOSE~$= 0$).
Sarle's bimodality coefficient
\(\BC = (\hat s^2 + 1) / (\hat\kappa + 3)\) is the bimodality
diagnostic (\(\BC > 0.555\) per Freeman \& Dale~\cite{Freeman2013}).
Results are summarised in \Cref{tab:t2}.

\begin{table}[h]
\centering
\caption{NET~2 results: bimodality coefficient ($\BC$), ATP statistics,
and commitment metrics across the nutrient axis (v9 model,
100 replicates per condition, run \texttt{20260704\_201120}).
$t_{\mathrm{depl}}$: nutrient depletion time ($\pm\,0$ min, deterministic).
Lag: time from depletion to ATP commitment dip.}
\label{tab:t2}
{\small\begin{tabular}{rcccrrr}
\toprule
$N_0$ ($\mu$M) & Spor.\,\% & $\BC$ & ATP\,floor (mM) & $t_{\mathrm{depl}}$ & Lag & Regime \\
\midrule
100  & 100 & 0.165 & 1.956 (spor) & 17 min & 140 min & fully committed \\
900  &  88 & 0.441 & 1.968 (spor) & 150 min & 127 min & near-commitment \\
1200 &  68 & 0.149 & 1.963 (spor) & 200 min & 117 min & bet-hedging \\
1440 &  38 & 0.356 & 1.970 (spor) & 240 min & 100 min & bet-hedging \\
     &     &       & 2.242 (veg)  &         &         & \\
1600 &   9 & 0.192 & 2.030 (spor) & 267 min &  87 min & tail-crossing \\
1800 &   0 & 0.270 & 3.428        & 300 min &  --- & vegetative \\
2000 &   0 & 0.533 & 4.435        & 333 min &  --- & vegetative \\
\bottomrule
\end{tabular}}
\end{table}

Final-ATP Sarle BC does \emph{not} exceed the $0.555$ bimodality
threshold in the commitment zone ($N_0 = 900$--$1600\,\mu$M).
The two ATP basins (sporulating: $\approx 1.96$~mM; vegetative:
$\approx 2.24$~mM at $N_0 = 1440$) are only $\approx 0.28$~mM
apart, and with $n = 100$ replicates the Sarle coefficient lacks the
sensitivity to separate two overlapping near-Gaussian peaks.
The composite order parameter $\varphi$ (eq.~\ref{eq:order-param})
combining Spo0A$\sim$P, $\sigma_H$, and Septum is required to recover
the true bimodal landscape at the commitment bifurcation.

\subsubsection*{Trajectory analysis: five unexpected landscape features}

G5-tier per-replicate trajectories reveal structure invisible in
the final-marking statistics:

\textbf{1. The commitment clock is deterministic.}
Nutrient depletion time has zero replicate-to-replicate variance
($\sigma = 0$~min) in all conditions: $t_{\mathrm{depl}} = N_0 /
(0.1\,\mu\mathrm{M\,s}^{-1}) / 60$.  The ODE-integrated nutrient
depletion sets an exact timer; only the post-depletion stochastic
phosphorelay introduces variability.

\textbf{2. Temporal priming: depletion lag decreases with $N_0$.}
The time between nutrient depletion and the ATP commitment dip
decreases from $140$~min at $N_0 = 100\,\mu$M to $87$~min at
$N_0 = 1600\,\mu$M (\Cref{tab:t2}).  Cells that deplete nutrients
later have been accumulating Spo0A and Spo0F longer; their phosphorelay
pools are larger when starvation hits, and they commit faster.
Spo0A$\sim$P peak confirms this: $13.8\,\mu$M at $N_0 = 100\,\mu$M
falls to $7.9\,\mu$M at $N_0 = 1600\,\mu$M.
This \emph{temporal priming effect} is invisible at the final-marking
level and can only be detected from trajectory data.

\textbf{3. The true commitment floor is 1.96 mM (trajectory minimum).}
The ATP dip during the commitment event is $1956$--$1970\,\mu$M
($= 1.96$~mM) regardless of $N_0$ --- a topology-derived constant.
This is distinct from the \emph{vegetative-basin floor} measured from
final markings ($\approx 2.24$~mM): cells whose ATP trajectory crosses
below $1.96$~mM commit; cells that do not cross recover to the
vegetative basin.  The $2.13$~mM value cited in earlier model versions
is an intermediate floor from a different parameter configuration;
the v9 value is $1.96$~mM (sporulation dip) and $2.24$~mM (vegetative
floor), a $0.28$~mM separation that defines the two landscape basins.

\textbf{4. Sporulation fraction = $P$($\sigma_H$ peak $> \theta_{\sigma H}$).}
At $N_0 = 1600\,\mu$M, mean $[\sigma_H]_{\max} = 1.41\,\mu$M
(below $\theta_{\sigma H} = 1.60\,\mu$M) yet $9\%$ sporulate.
The $9\%$ are exactly the replicates whose stochastic $\sigma_H$
trajectory transiently exceeds $1.60\,\mu$M.  The commitment fraction
is therefore a \emph{noise-driven threshold-crossing probability},
not a deterministic mean-field switch.

\textbf{5. SinR fluctuates in a narrow near-threshold band.}
Minimum SinR across all sporulating cells: $7.28$--$7.45\,\mu$M ---
only $0.05$--$0.22\,\mu$M below the $7.5\,\mu$M inhibitor threshold.
This near-miss geometry maximises the sensitivity of the SinI$\,/\,$SinR
gate to stochastic SinI bursts, amplifying small fluctuations into
binary fate decisions.

The landscape reveals a finding connecting directly to the hidden
variable in Fujita's experimental design.
The model predicts a \emph{commitment ATP floor} of $1.96$~mM
(trajectory minimum at the sporulation dip) and a
\emph{vegetative basin floor} of $\approx 2.24$~mM.
In \textbf{SM medium} (no carbon, cells nutrient-depleted) ATP declines
through the $1.96$~mM floor, enabling commitment.
In \textbf{CH medium} (rich carbon, exponential growth) ATP remains
above $2.24$~mM throughout, and the Spo0A$\sim$P pulse is cleared
before $[\sigma_H]$ builds.
The medium difference is therefore the \emph{mechanistic cause} of the
$52\%/5\%$ asymmetry, not a technical artefact of the inducible-promoter
strategy.  A direct falsifiable prediction: abrupt Spo0A* induction in
SM medium should achieve substantially higher sporulation than the
$\approx 5\%$ observed in CH.

\subsection{NET~3 --- Fluctuation-Theorem Inapplicability}
\label{ssec:t3}

A $2\times2$ factorial sweep was constructed using the
\texttt{INITIAL\_NUTRIENTS} $\times$ \texttt{NUTRIENTS\_STEP\_TARGET} axes
with a step event at $t = 3600$~s activated via \texttt{NUTRIENTS\_STEP\_TIME\_S}
$= 3600$~s as a fixed override.
Forward (\(N_0{=}1440 \to N_1{=}2160\)) and
reverse (\(N_0{=}2160 \to N_1{=}1440\)) protocols are compared with an
unperturbed baseline and two step-to-same controls (200 replicates per condition,
run \texttt{20260706\_151050}, v9 model).
Dissipated work $W$ was computed by the trajectory ATP-delta method
($W = [ATP_{\mathrm{init}} - ATP_{\mathrm{final}}] \times N_A \times V \times \Delta G_{\mathrm{ATP}}$).
Results are given in \Cref{tab:t3}.

\begin{table}[h]
\centering
\caption{NET~3 protocol-work distributions (v9 model, 200~replicates per condition,
$N_0 \in \{1440, 2160\}\,\mu$M step protocol).
$W_{\mathrm{fwd}}$ is negative because enrichment produces net ATP synthesis.}
\label{tab:t3}
\begin{tabular}{lrrrr}
\toprule
Condition & Spor.\,\% & $\overline{\mathrm{ATP}}$ (mM) & $\bar W / \kB T$ & Nuts$_{\mathrm{final}}$ ($\mu$M) \\
\midrule
Baseline ($N_0{=}1440$, no step) & 38 & 2.206 & $+3.72\times10^{7}$ & 0 \\
FWD ($1440{\to}2160$ at 1h)      &  0 & 5.201 & $-2.67\times10^{6}$ & 360 \\
REV ($2160{\to}1440$ at 1h)      &  2 & 3.398 & $+2.13\times10^{7}$ & 0 \\
Ctrl ($1440{\to}1440$ at 1h)     &  2 & 3.398 & $+2.13\times10^{7}$ & 0 \\
Ctrl ($2160{\to}2160$ at 1h)     &  0 & 5.206 & $-2.75\times10^{6}$ & 360 \\
\bottomrule
\end{tabular}
\end{table}

Three findings confirm strong irreversibility:

\textbf{1.  FT inapplicability.}
$|W / \kB T| = 3.72\times10^7$, equivalent to $10^{16\,000\,000}$ times the
Landauer bound.  The Jarzynski exponential average
$\langle e^{-W/\kB T}\rangle$ underflows in 64-bit floating point
($e^{-W/\kB T} < 10^{-10^7}$), and the Bennett-acceptance-ratio overlap
condition fails because forward and reverse distributions have no numerical
overlap.  This is not a numerical artefact; it is the direct signature of
POSet-enforced irreversibility.

\textbf{2.  FWD/REV work sign reversal (unexpected).}
$W_{\mathrm{fwd}} = -2.67\times10^6\,\kB T$ (negative: net ATP
\emph{production} when nutrients are enriched), while
$W_{\mathrm{rev}} = +2.13\times10^7\,\kB T$ (positive: ATP consumed).
The sign reversal means the forward and reverse protocols are not
time-reverses of each other --- a direct manifestation of the
asymmetric non-equilibrium landscape.

\textbf{3.  Biological irreversibility.}
Nutrient enrichment at 1~h completely aborts sporulation
(38\%~$\to$~0\%); nutrient depletion at 1~h only weakly triggers it
(0\%~$\to$~2\%).  Commitment is easy to prevent but nearly impossible
to induce by external perturbation alone.  The T3 negative result is the
strongest numerical confirmation that the POSet causal ordering enforces
the required irreversibility.



\subsection{NET~4 --- Critical Exponents at the Commitment Bifurcation}
\label{ssec:t4}

Two sweeps were combined: a wide grid ($N_0 \in \{1200,\ldots,2200\}\,\mu$M,
14~conditions, 300~replicates/condition, run~\texttt{20260706\_190903}) and a
dense sub-critical supplement centred on $N_c$ ($N_0 \in \{1250,\ldots,1360\}\,
\mu$M, 6~additional conditions, 300~replicates/condition,
run~\texttt{20260707\_134537}).
Results are given in \Cref{tab:t4}.

\begin{table}[h]
\centering
\caption{NET~4 combined results: sporulation fraction and Bernoulli $\CV^2$ of the
binary commitment outcome (300~reps/condition, v9 model).
Dense sub-critical points shaded.  $N_c = 1346\,\mu$M (linear interpolation at
the $50\%$ crossing).}
\label{tab:t4}
{\small\begin{tabular}{rrrrl}
\toprule
$N_0$ ($\mu$M) & Spor.\,\% & $\CV^2_{\mathrm{bin}}$ & $|N_0-N_c|$ ($\mu$M) & Side \\
\midrule
1200 & 69 & 0.449 & 146 & sub \\
\rowcolor{gray!15} 1250 & 64 & 0.562 &  96 & sub \\
\rowcolor{gray!15} 1280 & 60 & 0.657 &  66 & sub \\
1300 & 57 & 0.765 &  46 & sub \\
\rowcolor{gray!15} 1310 & 56 & 0.796 &  36 & sub \\
\rowcolor{gray!15} 1330 & 52 & 0.935 &  16 & sub \\
\rowcolor{gray!15} \textbf{1345} & \textbf{50} & \textbf{0.987} & \textbf{1} & sub \\
\rowcolor{gray!15} 1360 & 47 & 1.143 &  14 & sup \\
1380 & 44 & 1.256 &  34 & sup \\
1440 & 40 & 1.521 &  94 & sup \\
1500 & 33 & 2.030 & 154 & sup \\
1550 & 25 & 3.054 & 204 & sup \\
1600 & 14 & 6.143 & 254 & sup \\
\bottomrule
\end{tabular}}
\end{table}

The critical point is identified at $N_c = 1346\,\mu$M (dense grid interpolation
at $50\%$: $1345\,\mu$M gives $50.3\%$, $1360\,\mu$M gives $46.7\%$).
Power-law fits of $\CV^2_{\mathrm{bin}} \sim |N_0 - N_c|^{-\gamma}$ give
\begin{equation}
  \gamma_{\mathrm{sub}} = 0.150 \pm 0.044,\quad R^2 = 0.694,\quad n=7;
  \label{eq:gamma-v9}
\end{equation}
$\gamma_{\mathrm{sub}}$ excludes mean-field (\(\gamma = 1\)) at $19\,\sigma$.

The super-critical $\CV^2_{\mathrm{bin}}$ is dominated by Bernoulli statistics
($\CV^2 = (1-p)/p \to \infty$ as $p \to 0$) and does not yield a physical
exponent; $\gamma_{\mathrm{sup}}$ is not reported.

\textbf{Comparison with v7 model and interpretation.}
The v7 model (predecessor, $N_c = 21\,\mu$M) gave $\gamma_{\mathrm{sub}} = 2.44 \pm 0.33$
($R^2 = 0.9997$), a large super-mean-field exponent indicating strongly
fluctuation-dominated criticality.  The v9 value ($\gamma_{\mathrm{sub}} = 0.150$)
is much smaller, reflecting the wider commitment transition in v9 (spanning
$\sim 500\,\mu$M vs $\sim 10\,\mu$M in v7): the sporulation fraction changes
slowly with $N_0$, so the binary $\CV^2$ divergence is weak.
Both results share the key feature that $\gamma \neq 1$ (mean-field is excluded),
confirming that the commitment bifurcation is noise-driven rather than
gradient-driven in both model versions.

Trajectory analysis (G5 tier, dense sub-critical sweep
\texttt{run\_20260707\_134537}) reveals four additional features:

\textbf{1. Threshold-crossing mechanism.}
The mean $\sigma_H$ peak is $1.495$--$1.516\,\mu$M across all sub-critical
conditions $N_0 = 1250$--$1345\,\mu$M --- uniformly \emph{below} the
architectural threshold $\theta_{\sigma H} = 1.60\,\mu$M.
The sporulation fraction equals the probability that a stochastic $\sigma_H$
trajectory transiently exceeds $\theta_{\sigma H}$:
$P_{\mathrm{spor}} = P(\sigma_H^{\mathrm{stoch}} > 1.60\,\mu\mathrm{M})$.
As $N_0$ increases, the mean $\sigma_H$ peak decreases ($-2.4\times10^{-4}\,
\mu$M per $\mu$M nutrient), reducing the right-tail probability and lowering
the sporulation fraction smoothly.

\textbf{2. Exclusive Mode~B commitment.}
Every sporulating replicate at every sub-critical condition uses Mode~B
(SinI-mediated SinR sequestration at $t \approx 317$--$330\,\min$); zero
Mode~A events (passive early SinR degradation) were observed.
The bet-hedging mechanism near $N_c$ is therefore exclusively the stochastic
SinI burst --- the Mode~A pathway (which operates for shorter $\sigma_H$
half-lives, see NET~1) is inactive at the wild-type HL~$= 120\,\min$.

\textbf{3. Critical slowing down, modest magnitude.}
Mean commitment time increases from $317\,\min$ at $N_0 = 1250\,\mu$M (96~µM
below $N_c$) to $330\,\min$ at $N_0 = 1360\,\mu$M (14~µM above $N_c$): a
$+13\,\min$ ($4\%$) increase over a $110\,\mu$M approach.  This is consistent
with the weak $\gamma_{\mathrm{sub}} = 0.15$ --- the commitment clock is
dominated by the deterministic nutrient-depletion timer ($t_{\mathrm{depl}}
= N_0 / 0.1\,\mu\mathrm{M\,s}^{-1}$), with only a small stochastic add-on
at the bifurcation edge.

\textbf{4. Structural cascade lock-in at $N_c$.}
The covariance matrix at $N_0 = 1345\,\mu$M shows
$r(\sigma_G, \mathrm{Forespore}) = 0.996$, indicating that once the stochastic
SinI/SinR gate is crossed (the only stochastic step), the late execution
programme ($\sigma_G$ expression, forespore formation, cortex synthesis) is
essentially deterministic.  Commitment is the sole noise source; all downstream
structural transitions execute with near-perfect correlation.



\subsection{NET~5 --- Schnakenberg Entropy Production Rate Decomposition}
\label{ssec:t5}

Per-transition firing counts accumulated during simulation yield mean
firing rates \(\bar J_\rho = \bar n_\rho / t_{\mathrm{sim}}\), from
which the Schnakenberg EPR \eqref{eq:schnakenberg} was computed using
the NTP stoichiometry of each transition (\Cref{tab:t5-stoich});
results are reported in \Cref{tab:t5}.

\begin{table}[h]
\centering
\caption{NTP stoichiometry of stochastic transitions and per-firing free
  energy at \(T = 310.15\,\mathrm{K}\)
  (\(\Delta G_{\mathrm{ATP}} = \Delta G_{\mathrm{GTP}} = 23.27\,\kB T\)).}
\label{tab:t5-stoich}
\begin{tabular}{lcrl}
\toprule
Transition & ATP/firing & GTP/firing & \(\Delta G_\rho\) (\(\kB T\)/firing) \\
\midrule
\texttt{T\_KinA\_activation}      &  5 &  0 & 116.3 \\
\texttt{T\_Spo0F\_phosphorylation}&  0 &  0 & 0 (phosphotransfer) \\
\texttt{T\_Spo0F\_dephos}         &  0 &  0 & 0 (phosphatase) \\
\texttt{T\_Spo0A\_phosphorylation}&  0 &  0 & 0 (phosphotransfer) \\
\texttt{T\_Spo0A\_dephosphorylation}& 0 &  0 & 0 (phosphatase) \\
\texttt{T\_septation}             & 32 & 50 & 1\,908 \\
\bottomrule
\end{tabular}
\end{table}

\begin{table}[h]
\centering
\caption{NET~5 Schnakenberg EPR at three nutrient levels
  (300 replicates per condition, 6\,h horizon).
  \(\dot\sigma_{\mathrm{total}}\) is the sum of per-transition contributions.
  Phosphorelay transitions contribute \(\dot\sigma = 0\) by stoichiometry.}
\label{tab:t5}
\begin{tabular}{rrrrrr}
\toprule
\(N\) (\(\mu\)M) & \(\bar n_{\mathrm{KinA}}\) & \(\bar n_{\mathrm{sept}}\)
  & \(\dot\sigma_{\mathrm{KinA}}\) (\(\kB T/\mathrm{s}\))
  & \(\dot\sigma_{\mathrm{sept}}\) (\(\kB T/\mathrm{s}\))
  & \(\dot\sigma_{\mathrm{total}}\) (\(\kB T/\mathrm{s}\)) \\
\midrule
15  & 88.7  & 144.3 & 288 & 7\,675  & 7\,963  \\
\textbf{21}  & 73.2  & 182.2 & 237 & 9\,690  & 9\,928  \\
30  & 78.8  & 192.5 & 255 & 10\,240 & 10\,496 \\
\bottomrule
\end{tabular}
\end{table}

Two findings:
(i)~\emph{Septation dominates dissipation} at every condition
(\(\geq 96\%\)).  The KinA autophosphorylation contributes \(\sim 3\%\);
the phosphorelay transitions themselves contribute zero explicit NTP
dissipation by stoichiometry.  The \emph{decision} costs \(< 4\%\) of the
circuit energy budget; the \emph{execution} of the chosen fate costs
\(\sim 96\%\).
(ii)~EPR is \emph{monotonically} ordered growth (7\,963) \(<\) critical
(9\,928) \(<\) starvation (10\,496) \(\kB T / \mathrm{s}\).  This is
inconsistent with a ``maximum entropy production at criticality''
hypothesis; the bifurcation does not coincide with peak dissipation.

At the critical point the absolute power dissipated by the circuit is
\(\sim 4.3\times10^{-17}\)~W, which is \(\sim 4\times10^{-4}\) of the
total metabolic budget of a growing \textit{B.~subtilis} cell
(\(\sim 10^{-13}\)~W~\cite{Noor2016}) --- physically plausible for a
single regulatory module.

\subsection{Summary: Five-Way Co-Validation and Fujita Correspondence}
\label{ssec:covalidation}

\Cref{tab:covalidation} summarises the bidirectional co-validation:
each NET test simultaneously confirms a thermodynamic property of the
open-system transition, a structural feature of the SHPN architecture,
\emph{and} provides a physical answer to one of the three Fujita
questions (Q1--Q3) that motivated the biological validation.

\begin{table}[h]
\centering
\caption{Bidirectional co-validation and Fujita correspondence.
  Each NET test confirms a thermodynamic property, an architectural feature,
  and provides a physics-level answer to one of the three
  Fujita~\cite{Fujita2005} questions.  The three columns are not independent:
  they are the same fact in three languages (physics, formalism, experiment).}
\label{tab:covalidation}
{\small
\begin{tabular}{p{1.8cm}p{3.8cm}p{3.5cm}p{3.5cm}}
\toprule
Test & Thermodynamic property & Architecture feature & Fujita question \\
\midrule
NET~1 (\(\eta_L\))
  & \(\eta_L \propto n\): sharper gates approach Landauer bound;
    monotonic improvement with half-life
  & \(\Gamma = (K,n,\varepsilon)\) is the experimentally grounded gate spec
  & \textbf{Q1 (partial)}: gradual route achieves higher \(\eta_L\) than
    abrupt pulse (more efficient Landauer step) \\
NET~2 (\(U_{\mathrm{eff}}\))
  & Two-basin landscape; \emph{vegetative basin ATP floor} at $2.13$~mM
    is a topology-driven prediction (no parameter fitting)
  & \(\Ge \times \Gs\): attractor count is a property of \(\Ge\) alone
  & \textbf{Q2}: ATP floor marks the thermodynamic boundary below which
    the vegetative programme is abandoned \\
NET~3 (FT)
  & \(W/\kB T \sim 10^8\): equilibrium FTs inapplicable
  & POSet \(\lambda\): syntactic causal ordering enforces irreversibility
  & \textbf{Q1 (partial)}: dissipated work scale rules out reversibility ---
    neither the gradual nor the abrupt route can uncommit \\
NET~4 (\(\gamma_\pm\))
  & \(\gamma_{\mathrm{sub}} \approx 2.4\), \(\gamma_{\mathrm{sup}} \approx 4.6 \gg 1\):
    noise-driven, not gradient-driven
  & \(\Gamma\)-gated \(\Gs\): topology, not rate magnitudes, sets the
    bifurcation class
  & \textbf{Q3}: non-mean-field exponents produce equal basin depths ---
    $\approx 50\%$ bet-hedging is the thermodynamically natural outcome
    for this universality class \\
NET~5 (Schnak.)
  & \(\dot\sigma_{\Gs}/\dot\sigma_{\mathrm{total}} < 4\%\);
    septation dominates (\(>96\%\))
  & \(\Ge \times \Gs\): orthogonal entropy flows; EPR decomposable by subgraph
  & \textbf{Q2}: $\sigma_H$/septation gate concentrates the entire
    irreversibility budget --- $\sigma_H$-null eliminates the thermodynamic
    commitment mechanism, not merely a downstream step \\
\bottomrule
\end{tabular}}
\end{table}

% ─────────────────────────────────────────────────────────────────────────────
\subsection{Fujita Protocol Validation: Three Biological Questions Answered}
\label{ssec:fujita_validation}

The three questions posed by Fujita and Losick~\cite{Fujita2005} were
addressed through four targeted parameter sweeps (F1--F4,
\Cref{tab:fujita_sweeps}) on the v9 SHPN model (runs
F1~=~\texttt{run\_20260702\_171823}, F2~=~\texttt{run\_20260702\_180944},
F3~=~\texttt{run\_20260702\_191849}).
Each sweep reports the sporulation fraction (Mature\_spore\_final$>0.5$
across replicates), the mean peak $\sigma_H$ concentration, and the
vegetative-basin ATP floor (mean trajectory minimum across non-sporulating
replicates).
Key results are summarised in \Cref{tab:fujita_results}; the full
condition-level data are archived in the repository.

\medskip\noindent\textbf{Q1 --- Timing and medium both matter, but the medium is the gate.}

Sweep~F1 showed that the sporulation fraction under identical nutrient
conditions depends critically on whether the Spo0A$\sim$P accumulation
is gradual (phosphorelay, \texttt{LOADING\_DOSE}~$= 0$) or abrupt
(single pulse, \texttt{LOADING\_DOSE}~$= 27\,\mu$M added at $t = 1$~s):
\begin{itemize}
\item \textbf{SM-like} (\texttt{INITIAL\_NUTRIENTS}~$= 100\,\mu$M, fast
  starvation): \emph{both} routes gave $100\%$ sporulation.  The abrupt pulse
  merely shortened mean commitment time from 151 to 136~min.
\item \textbf{CH-like} (\texttt{INITIAL\_NUTRIENTS}~$= 1440\,\mu$M,
  slow starvation): the gradual route gave $38\%$; the abrupt pulse gave
  $56\%$.  Neither achieved the $100\%$ seen in SM medium, and neither
  reproduced the Fujita $\approx5\%$ seen with constitutive Spo0A*
  in rich medium.
\end{itemize}

Sweep~F2 extended the result over the full nutrient axis with the abrupt
pulse fixed at $27\,\mu$M.  The sporulation fraction decreased
\emph{monotonically} from $100\%$ at $N_0 = 100\,\mu$M to $0\%$ at
$N_0 = 2200\,\mu$M (\Cref{tab:fujita_results}).  The $50\%$ crossover
fell between 1440 and $1600\,\mu$M, identifying the effective ATP
commitment boundary.  The nutrient depletion analysis shows why:
nutrients deplete at a constant rate of $0.1\,\mu$M~s$^{-1}$ (model
parameter), so only conditions with $N_0 \leq 2160\,\mu$M reach full
depletion within the 6-h simulation window; cells in richer medium
(\texttt{INITIAL\_NUTRIENTS}~$> 2160\,\mu$M) remain in an energy-replete
state throughout, and KinA activation --- gated by
$[1 + (\text{Nutrients})^2]^{-1}$ --- is negligible.

These findings together answer Q1 thermodynamically: the relevant
hidden variable between the Fujita SM and CH experiments is the
\emph{nutrient-depletion trajectory}, not merely the abruptness of
Spo0A induction.  SM medium allows the ATP pool to fall below the
vegetative-basin floor ($2.22 \pm 0.04$~mM, measured as the mean
trajectory minimum across non-sporulating replicates); CH medium keeps
it above that floor indefinitely.

\medskip\noindent\textbf{Q2 --- $\sigma_H$ is the thermodynamic gate, not a downstream reporter.}

Sweep~F3 used the new \texttt{k\_sigmaH\_factor} parameter (P41 in the
v9 model), which scales the $\sigma_H$ transcription rate, allowing a
clean titration from fully silenced ($= 0$, true $\sigma_H$-null) to
fully active ($= 1$, wild-type).  With \texttt{INITIAL\_NUTRIENTS}~$= 100$
(fast starvation) and either route ($\texttt{LOADING\_DOSE} \in \{0, 27\}$):
\begin{itemize}
\item $k_{\sigma H} = 0$ (true $\sigma_H$-null): $\mathbf{0\%}$
  sporulation with both natural phosphorelay \emph{and} abrupt pulse.
  Peak $[\sigma_H] = 0.000\,\mu$M.  This reproduces the Fujita
  Fig.~4 result.
\item $k_{\sigma H} = 0.2$: $96\%$ (natural) and $100\%$ (pulse).
  Even $20\%$ $\sigma_H$ transcription capacity is sufficient for
  near-complete commitment at SM nutrient levels --- the transition
  from null to active is essentially a step.
\item $k_{\sigma H} = 1.0$: $100\%$ for both routes, mean peak
  $[\sigma_H] \approx 1.64\,\mu$M (above $\theta_{\sigma H} = 1.60\,\mu$M$).
\end{itemize}
The sharp null/active threshold, combined with the fact that the
abrupt Spo0A$\sim$P pulse cannot rescue commitment in the absence of
$\sigma_H$, confirms that $\sigma_H$ is the \emph{irreversible gate}
in the cascade: no matter how much upstream signal is present, the cell
cannot commit without it.

\begin{table}[ht]
\centering
\caption{\textbf{Fujita Protocol sweep results} (v9 model, 6-h horizon,
  \texttt{INITIAL\_NUTRIENTS}~$= 100\,\mu$M fixed in F1/F3, swept in F2).
  Spor\,\%: sporulation fraction.  $[\sigma_H]_{\max}$: mean peak
  $\sigma_H$ across all replicates.
  ATP$_{\mathrm{veg}}$: mean ATP trajectory minimum across
  non-sporulating replicates (vegetative basin floor).
  $t_{\mathrm{commit}}$: mean time to first mature spore (sporulating
  replicates only).
  F4 results: run\_20260704\_163628, v9 model, server SHA 7057e706.}
\label{tab:fujita_results}
{\small
\begin{tabular}{@{}llccccc@{}}
\toprule
Sweep & Condition & Spor\,\% & $[\sigma_H]_{\max}$ ($\mu$M)
  & ATP$_{\mathrm{veg}}$ (mM) & $t_{\mathrm{commit}}$ (min) \\
\midrule
\multicolumn{6}{@{}l}{\textit{F1 --- medium $\times$ route factorial}} \\
& SM natural ($N_0{=}100$, D$=0$) & 100 & 1.635 & --- & 151 \\
& SM abrupt  ($N_0{=}100$, D$=27$) & 100 & 1.647 & --- & 136 \\
& CH natural ($N_0{=}1440$, D$=0$) & 38  & 1.464 & 2.228 & 335 \\
& CH abrupt  ($N_0{=}1440$, D$=27$) & 56  & 1.553 & 2.222 & 327 \\
\midrule
\multicolumn{6}{@{}l}{\textit{F2 --- N$_0$ titration, abrupt pulse fixed}} \\
& $N_0{=}100\,\mu$M   & 100 & 1.647 & --- & 136 \\
& $N_0{=}600\,\mu$M   & 100 & 1.636 & --- & 207 \\
& $N_0{=}1200\,\mu$M  &  81 & 1.595 & 2.223 & 300 \\
& $N_0{=}1440\,\mu$M  &  56 & 1.564 & 2.222 & 327 \\
& $N_0{=}1800\,\mu$M  &   1 & 1.452 & 3.367 & 332 \\
& $N_0{=}2200\,\mu$M  &   0 & --- & 5.000 & --- \\
\midrule
\multicolumn{6}{@{}l}{\textit{F3 --- $\sigma_H$ gate titration, SM medium, $2\times2$}} \\
& $k_{\sigma H}{=}0$, D$=0$ (null natural) & 0 & 0.000 & --- & --- \\
& $k_{\sigma H}{=}0$, D$=27$ (null + pulse) & 0 & 0.000 & --- & --- \\
& $k_{\sigma H}{=}0.2$, D$=0$ & 96 & 0.320 & --- & 182 \\
& $k_{\sigma H}{=}0.2$, D$=27$ & 100 & 0.327 & --- & 163 \\
& $k_{\sigma H}{=}1.0$, D$=0$ (wild-type) & 100 & 1.639 & --- & 147 \\
& $k_{\sigma H}{=}1.0$, D$=27$ & 100 & 1.648 & --- & 132 \\
\midrule
\multicolumn{6}{@{}l}{\textit{F4 --- natural route $N_0$ titration (200 reps/cond)}} \\
& $N_0{=}100\,\mu$M (commitment) & 100 & 1.569 & 2.166 & --- \\
& $N_0{=}1000\,\mu$M & 84 & 1.319 & 2.172 & --- \\
& $N_0{=}1300\,\mu$M & 56 & 1.179 & 2.190 & --- \\
& $N_0{=}1440\,\mu$M & 39 & 1.093 & 2.206 & --- \\
& $N_0{=}1800\,\mu$M & 2 & 0.840 & 3.398 & --- \\
& $N_0{=}2000\,\mu$M (vegetative) & 0 & 0.681 & 4.434 & --- \\
\multicolumn{3}{@{}l}{\textit{$N_c \approx 1370\,\mu$M; bet-hedging window 700--1800~$\mu$M}} & & & \\
\bottomrule
\end{tabular}}
\end{table}

% ─────────────────────────────────────────────────────────────────────────────
\section{Discussion}
\label{sec:discussion}

\subsection{The Two-Pillar Co-Validation Structure}

The contribution of this work is the demonstration that two independent
pillars --- the physics of non-equilibrium open-system thermodynamics, and
the graph-theoretic architecture of SHPN --- jointly characterise
biological commitment in a way that neither pillar can on its own.  The
direction \emph{architecture \(\to\) physics} is shown by NET~4 and NET~5:
the \(\Gamma\)-gated POSet topology of \(\Gs\) predicts non-mean-field
critical exponents and an asymmetric EPR decomposition that are then
measured numerically.  The direction \emph{physics \(\to\) architecture} is
shown by NET~3: the dissipated work \(W/\kB T \sim 10^8\) measured from
the simulation is the empirical signature confirming that the POSet causal
ordering is enforcing strong irreversibility, exactly as the formalism
requires.

A particularly informative finding is the NET~5 decomposition:
the decision step (\(\Gs\)) costs only a few percent of the circuit energy
budget, while the execution step (\(\Ge\): septation, coat synthesis)
costs the bulk.  This was not obvious \textit{a~priori} and is impossible
to read off any formalism that does not separate the two subgraphs.
The NET~2 landscape provides an external calibration:
the vegetative basin floor at $\approx 2.13$~mM is a
topology-driven prediction that requires no parameter fitting
to the commitment window;
this adds a third validation direction:
\emph{topology predicts a physically observable threshold}.

\subsection{Physics Answers to the Fujita Questions}
\label{ssec:fujita-answers}

The preceding Results section (\Cref{ssec:fujita_validation}) confirmed
the three Fujita questions with direct simulation evidence.  Here we
translate each answer into the thermodynamic vocabulary of the framework.

\textbf{Q1 answered: the medium sets the thermodynamic gate, the route is
secondary.}
Sweep~F1 showed that SM medium ($N_0 = 100\,\mu$M) gives $100\%$
sporulation with \emph{either} route, while CH medium ($N_0 = 1440\,\mu$M)
gives $38$--$56\%$ regardless of whether a pulse is applied.
Sweep~F2 confirmed the monotonic dependence: sporulation decreases
strictly with $N_0$, falling to $0\%$ at $N_0 = 2200\,\mu$M.
The physical explanation is the vegetative-basin ATP floor.  In SM
medium nutrients deplete at $0.1\,\mu$M~s$^{-1}$, driving the ATP pool
below $2.22$~mM within $\sim 17$~min; this depletion activates KinA
(rate $\propto [1 + N^2]^{-1}$, near-maximal only when $N \approx 0$)
and initiates the phosphorelay.  In CH medium the cell remains
energy-replete throughout the 6-h window --- the ATP floor is never
reached, KinA activation is negligible, and the Spo0A$\sim$P pulse
(whether from the natural relay or an abrupt injection) is cleared by
the nutrient-boosted phosphatase before $\sigma_H$ can accumulate.
In Landauer terms: the abrupt route in rich medium cannot complete the
free-energy storage step ($\sigma_H$ ramp $\to$ septation) because the
energy substrate (ATP depletion below the floor) is not available.
The $5\%$ Fujita result arises from the rare stochastic escape above
the ATP floor even at $N_0 = 2200\,\mu$M; our model gives $0\%$ at
that condition, consistent with the experimental bound within sampling
uncertainty (Fujita used $< 200$~cells per time point).

\textbf{Q2 answered: $\sigma_H$ is the circuit's Landauer step.}
Sweep~F3 used the \texttt{k\_sigmaH\_factor} parameter to titrate
$\sigma_H$ transcription from fully active ($= 1$, wild-type) to
completely silenced ($= 0$, true $\sigma_H$-null).  At $k_{\sigma H} = 0$:
$0\%$ sporulation with both natural and abrupt route, confirming
Fujita Fig.~4.  The on/off threshold is extremely sharp: even
$k_{\sigma H} = 0.2$ restores $96$--$100\%$ commitment.  The physics
interpretation follows directly from NET~5 (Schnakenberg decomposition,
v7 model): septation, the transition gated by $\sigma_H$, concentrates
$> 96\%$ of total circuit entropy production ($\dot\sigma_{\mathrm{sept}} =
9928\,\kB T\,\text{s}^{-1}$ at the critical condition vs.\ $237\,\kB T\,\text{s}^{-1}$
for KinA autophosphorylation).  $\sigma_H$-null therefore does not merely
remove a downstream reporter; it removes the \emph{only} step where the
circuit pays its Landauer cost.  Without septation, the cell cannot
execute the irreversible state transition, regardless of how much
upstream Spo0A$\sim$P is present.

\textbf{Q3 answered: $N_c \approx 1370\,\mu$M and the bet-hedging window
spans 700--1800~$\mu$M.}
Sweep~F4 (FUJITA-4, 200 replicates $\times$ 16 conditions, natural route,
\texttt{run\_20260704\_163628}) resolved the full sporulation fraction curve
on the v9 model at single-condition resolution:
\begin{itemize}
\item \textbf{Commitment zone} ($N_0 \leq 400\,\mu$M): $100\%$
  sporulation (95\% CI [98--100\%]).  Nutrients deplete fully;
  $[\sigma_H]_{\max} \approx 1.52$--$1.57\,\mu$M $>\theta_{\sigma H}$.
\item \textbf{Bet-hedging window} (700--1800~$\mu$M): sporulation falls
  monotonically from $96\%$ [92--98\%] at $N_0 = 700\,\mu$M to
  $2\%$ [1--4\%] at $N_0 = 1800\,\mu$M.  At $N_0 = 1300\,\mu$M:
  $56\%$ [49--62\%]; at $N_0 = 1440\,\mu$M: $39\%$ [32--45\%].
  The $50\%$ crossover is $N_c \approx 1370\,\mu$M (midpoint of the
  first condition above and the first condition below 50\%;  95\%~CIs
  are non-overlapping, confirming the crossing is statistically resolved).
\item \textbf{Vegetative zone} ($N_0 \geq 2000\,\mu$M): $0\%$
  sporulation [0--2\%]; nutrients are not fully depleted within 6~h;
  $[\sigma_H]_{\max}$ falls to $0.68\,\mu$M at $N_0 = 2000\,\mu$M,
  well below $\theta_{\sigma H} = 1.60\,\mu$M.
\end{itemize}
The vegetative ATP floor is $2.165$--$2.225$~mM across all
nutrient-stressed conditions ($N_0 \leq 1800\,\mu$M), consistent with
the topology-derived prediction of $2.13$~mM and with the F1/F2
measurements ($2.22 \pm 0.04$~mM).  The NET~4 critical exponents
($\gamma_{\mathrm{sub}} \approx 2.4$, $\gamma_{\mathrm{sup}} \approx 4.6$,
v7 model) predict asymmetric basin curvature that naturally produces a
near-50\% crossover near the bifurcation; the F4 measurement of
$N_c \approx 1370\,\mu$M confirms that the Fujita $\approx52\%$ is the
thermodynamically natural outcome near the noise-driven commitment
bifurcation, not a biologically tuned value.

\subsection{Competitive Landscape}

SHPN's competitors are biological modelling formalisms, not theoretical
physics.  \Cref{tab:competition} compares SHPN to the principal
alternatives on the five criteria required for the co-validation of
\Cref{tab:covalidation}.

\begin{table}[h]
\centering
\caption{Comparison of biological modelling formalisms on the five
criteria required for thermodynamic characterisation of commitment.
SHPN is the only formalism satisfying all five.}
\label{tab:competition}
\begin{tabular}{lccccc}
\toprule
Formalism & Stoch.\ & Signal & Causal & Thermodyn. & Hybrid \\
          & dynamics & hierarchy & ordering & \(\Gamma\) & stoch/ODE \\
\midrule
SBML / ODEs                & \checkmark & --- & --- & --- & partial \\
Standard SPN / GSPN        & \checkmark & --- & --- & --- & --- \\
Boolean / Thomas networks  & qualitative & --- & --- & --- & --- \\
Kappa / BioNetGen          & \checkmark & --- & --- & --- & --- \\
Harel Statecharts          & partial & \checkmark & \checkmark & --- & --- \\
Bio-CCS / process algebras & \checkmark & --- & --- & --- & --- \\
\textbf{SHPN}              & \checkmark & \checkmark & \checkmark & \checkmark & \checkmark \\
\bottomrule
\end{tabular}
\end{table}

The closest match on the structural criteria (signal hierarchy + causal
ordering) is Harel Statecharts~\cite{Harel1987}, but Statecharts lack
stoichiometry, stochastic kinetics, and any thermodynamic interpretation
of their guards.  SHPN is, to our knowledge, the first biological
modelling formalism in which each architectural decision carries an
explicit thermodynamic meaning testable by simulation.

\subsection{Limitations and Open Questions}

\textit{Tightness of the Landauer bound.}
\(\eta_L \sim 10^{-9}\) measured by NET~1 establishes the bound is far
from saturated by the SHPN model.  Whether the gap is intrinsic to the
biology, an artefact of model coarse-graining, or an indication of
overestimated structural-execution costs requires independent
calorimetric measurement.

\medskip\noindent\textit{Universality class identification.}
NET~4 excludes mean-field but cannot identify a universality class on the
basis of two power-law exponents alone.  Finite-size scaling analysis
across a range of cell-volume scales would be required for a rigorous
classification.

\medskip\noindent\textit{Multi-layer PreemptionCheck.}
The current implementation enforces PreemptionCheck on direct
signal predecessors only (single-layer, non-recursive).  Multi-layer
enforcement requires a layer-assignment computation on \(\Gs\) that is
not yet integrated into the runtime.  Whether the single-layer
approximation under- or over-estimates the Landauer cost of the cascade
is an open question.

\medskip\noindent\textit{Beyond binary commitment.}
\(|\Aset(\Ge)| = 2\) is the simplest non-trivial case.  Cells with
larger fate spaces (e.g.\ haematopoietic stem cells with \(\sim 10\)
distinguishable lineages) are the natural next test.  The
\(\kB T \ln|\Aset(\Ge)|\) scaling predicted by the framework is
falsifiable by such systems.

\subsection{Positioning}

SHPN is a domain-specific formal language for biological
decision-making with an intrinsic open-system thermodynamic semantics.
It extends stochastic Petri nets with three architectural primitives ---
\(\Ge \times \Gs\), POSet \(\lambda\) with PreemptionCheck, and the
\(\Gamma\) gate tuple --- that together formalise the three quantifiable
components of the Waddington landscape:
\begin{enumerate}
\item the attractor space \(\Aset(\Ge)\), with minimum selection cost
  \(\kB T \ln|\Aset(\Ge)|\) derivable from topology;
\item the landscape tilting mechanism: \(\Gs\) driven by
  \(\Gamma\)-gated arcs whose parameters \((K, n, \varepsilon)\) are
  grounded in dose-response measurements;
\item the entropy cost of fate selection, decomposable as
  \(\dot\sigma = \dot\sigma_{\Ge} + \dot\sigma_{\Gs}\) via Schnakenberg.
\end{enumerate}
This positions the work in the tradition of Ferrell~\cite{Ferrell1998},
Tyson and Novak~\cite{TysonNovak2001}, Thomas~\cite{Thomas1973},
and S\"uel et~al.~\cite{Suel2007} --- formalisms answering \emph{how cells
decide} --- while being the first to make signal hierarchy and its
causal ordering first-class syntactic features and to assign each
architectural primitive an explicit thermodynamic meaning.

% ─────────────────────────────────────────────────────────────────────────────
\section{Conclusion}
\label{sec:conclusion}

We have argued that cell-fate commitment is best described as an
irreversible open-system thermodynamic transition, and that the
graph-theoretic architecture of Signal Hierarchical Petri Nets (SHPN) is the formal
modelling language that makes the thermodynamic claim computable,
testable, and biologically grounded.  Five numerical experiments on a
SHPN model of \textit{B.~subtilis} sporulation jointly establish a
bidirectional co-validation between the framework's architecture and its
thermodynamic predictions: the model's Landauer efficiency
(\(\sim 10^{-9}\)) and its monotonic improvement with switch sharpness,
the two-basin non-equilibrium landscape, the dissipated-work magnitude
(\(\sim 10^{8}\,\kB T\)), the non-mean-field critical exponents at the
nutrient bifurcation, and the dominance of structural execution over
informational decision in the Schnakenberg EPR decomposition.  No prior
biological modelling formalism allows this characterisation to be
performed end-to-end within the formalism.  Validated here on a binary
commitment system, the framework is applicable to any cell-fate
circuit for which the regulatory topology is known; extension to larger
fate repertoires --- haematopoietic lineage decisions, neural subtype
specification --- where the \(\kB T \ln|\Aset(\Ge)|\) scaling is
directly falsifiable, is a natural next step.

% ─────────────────────────────────────────────────────────────────────────────
\section*{Author Contributions}
\addcontentsline{toc}{section}{Author Contributions}

E.S.\ conceived the study, developed the SHPN formalism and the SHyPN
software platform, designed and dispatched the sweeps, analysed the
results, and wrote the manuscript.

% ─────────────────────────────────────────────────────────────────────────────
\section*{Data Availability}
\addcontentsline{toc}{section}{Data Availability}

\begin{sloppypar}
All materials are publicly available at
\url{https://github.com/simao-eugenio/shypn}
(DOI:~\href{https://doi.org/10.5281/zenodo.20835073}{10.5281/zenodo.20835073})
under the MIT open source licence;
full details of model files, sweep artefacts, and analysis scripts
are given in the Reproducibility subsection of Methods.
Statistical analyses were performed in Python~3.11 using
\mbox{pandas~2.x}, NumPy~1.26.x, and SciPy~1.12.x.
No proprietary software was used.
\end{sloppypar}

% ─────────────────────────────────────────────────────────────────────────────
\section*{Use of Artificial Intelligence Tools}
\addcontentsline{toc}{section}{Use of Artificial Intelligence Tools}

During the preparation of this manuscript, the author used
GitHub Copilot (powered by large language models) to assist with code
generation (simulation scripts, analysis pipelines, figure scripts) and
manuscript proofreading.
All scientific content, theoretical development, experimental design,
data analysis, and conclusions are the sole responsibility of the author.
No AI tool was used to generate original scientific claims or to replace
critical evaluation.

% ─────────────────────────────────────────────────────────────────────────────
\section*{Acknowledgments}
\addcontentsline{toc}{section}{Acknowledgments}

This work was performed within the SHyPN project at the Federal University
of Santa Catarina (UFSC), Brazil.
Computational resources were provided by a dedicated GPU server
(NVIDIA RTX 5060~Ti, CUDA~12.9).
The author thanks the Fujita and Losick laboratories for the published
experimental measurements that calibrated the \textit{B.~subtilis}
sporulation model.

% ─────────────────────────────────────────────────────────────────────────────
\begin{thebibliography}{99}

\bibitem{Waddington1957}
C.~H. Waddington,
\textit{The Strategy of the Genes}.
George Allen \& Unwin, 1957.

\bibitem{Ferrell1998}
J.~E. Ferrell~Jr.\ and E.~M. Machleder,
``The biochemical basis of an all-or-none cell fate switch in
\textit{Xenopus} oocytes,''
\textit{Science}, vol.~280, pp.~895--898, 1998.

\bibitem{TysonNovak2001}
J.~J. Tyson, K.~C. Chen, and B.~Novak,
``Network dynamics and cell physiology,''
\textit{Nature Reviews Molecular Cell Biology}, vol.~2, pp.~908--916, 2001.

\bibitem{Thomas1973}
R.~Thomas,
``Boolean formalisation of genetic control circuits,''
\textit{Journal of Theoretical Biology}, vol.~42, pp.~563--585, 1973.

\bibitem{Suel2007}
G.~M. S\"{u}el, R.~P. Kulkarni, J.~Dworkin, J.~Garcia-Ojalvo, and
M.~B. Elowitz,
``Tunability and noise dependence in differentiation dynamics,''
\textit{Science}, vol.~315, pp.~1716--1719, 2007.

\bibitem{Wang2011}
J.~Wang, K.~Zhang, L.~Xu, and E.~Wang,
``Quantifying the Waddington landscape and biological paths for
development and differentiation,''
\textit{Proceedings of the National Academy of Sciences}, vol.~108,
pp.~8257--8262, 2011.

\bibitem{Nicolis1977}
G.~Nicolis and I.~Prigogine,
\textit{Self-Organization in Nonequilibrium Systems}.
Wiley-Interscience, 1977.

\bibitem{Landauer1961}
R.~Landauer,
``Irreversibility and heat generation in the computing process,''
\textit{IBM Journal of Research and Development}, vol.~5, pp.~183--191, 1961.

\bibitem{Schnakenberg1976}
J.~Schnakenberg,
``Network theory of microscopic and macroscopic behavior of master
equation systems,''
\textit{Reviews of Modern Physics}, vol.~48, pp.~571--585, 1976.

\bibitem{Jarzynski1997}
C.~Jarzynski,
``Nonequilibrium equality for free energy differences,''
\textit{Physical Review Letters}, vol.~78, pp.~2690--2693, 1997.

\bibitem{Crooks1999}
G.~E. Crooks,
``Entropy production fluctuation theorem and the nonequilibrium work
relation for free energy differences,''
\textit{Physical Review E}, vol.~60, pp.~2721--2726, 1999.

\bibitem{Marcus1993}
R.~A. Marcus,
``Electron transfer reactions in chemistry: theory and experiment,''
\textit{Reviews of Modern Physics}, vol.~65, pp.~599--610, 1993.

\bibitem{Marsan1995}
M.~Ajmone~Marsan, G.~Balbo, G.~Conte, S.~Donatelli, and G.~Franceschinis,
\textit{Modelling with Generalized Stochastic Petri Nets}.
Wiley, 1995.

\bibitem{Wilkinson2018}
D.~J. Wilkinson,
\textit{Stochastic Modelling for Systems Biology},
3rd ed.
Chapman \& Hall/CRC, 2018.

\bibitem{Harel1987}
D.~Harel,
``Statecharts: a visual formalism for complex systems,''
\textit{Science of Computer Programming}, vol.~8, pp.~231--274, 1987.

\bibitem{Freeman2013}
J.~B. Freeman and J.~Dale,
``Assessing bimodality to detect the presence of a dual cognitive
process,''
\textit{Behavior Research Methods}, vol.~45, pp.~83--97, 2013.

\bibitem{Noor2016}
E.~Noor \textit{et al.},
``The protein cost of metabolic fluxes: prediction from enzymatic rate
laws and cost minimization,''
\textit{PLOS Computational Biology}, vol.~12, e1005167, 2016.

\bibitem{MehtaSchwab2012}
P.~Mehta and D.~J. Schwab,
``Energetic costs of cellular computation,''
\textit{Proceedings of the National Academy of Sciences}, vol.~109,
pp.~17978--17982, 2012.

\bibitem{Berut2012}
A.~B\'erut, A.~Arakelyan, A.~Petrosyan, S.~Ciliberto, R.~Dillenschneider,
and E.~Lutz,
``Experimental verification of Landauer's principle linking information
and thermodynamics,''
\textit{Nature}, vol.~483, pp.~187--189, 2012.

\bibitem{Fujita2005}
M.~Fujita and R.~Losick,
``Evidence that entry into sporulation in \textit{Bacillus subtilis} is
governed by a gradual increase in the level and activity of the master
regulator {Spo0A},''
\textit{Genes \& Development}, vol.~19, pp.~2236--2244, 2005.

\bibitem{Narula2012}
J.~Narula, S.~N. Devi, M.~Fujita, and O.~A. Igoshin,
``Ultrasensitivity of the \textit{Bacillus subtilis} sporulation
decision,''
\textit{Proceedings of the National Academy of Sciences}, vol.~109,
pp.~E3513--E3522, 2012.

\bibitem{Burbulys1991}
D.~Burbulys, K.~A. Trach, and J.~A. Hoch,
``Initiation of sporulation in \textit{B.~subtilis} is controlled by a
multicomponent phosphorelay,''
\textit{Cell}, vol.~64, pp.~545--552, 1991.

\bibitem{Stragier1996}
P.~Stragier and R.~Losick,
``Molecular genetics of sporulation in \textit{Bacillus subtilis},''
\textit{Annual Review of Genetics}, vol.~30, pp.~297--341, 1996.

\bibitem{Errington2003}
J.~Errington,
``Regulation of endospore formation in \textit{Bacillus subtilis},''
\textit{Nature Reviews Microbiology}, vol.~1, pp.~117--126, 2003.

\bibitem{BRENDA}
I.~Schomburg \textit{et al.},
``{BRENDA} in 2013: integrated reactions, kinetic data, enzyme function
data, improved disease classification,''
\textit{Nucleic Acids Research}, vol.~41, pp.~D764--D772, 2013.

\bibitem{GeQian2010}
H.~Ge and H.~Qian,
``Physical origins of entropy production, free energy dissipation, and
their mathematical representations,''
\textit{Physical Review E}, vol.~81, p.~051133, 2010.

\bibitem{Ao2004}
P.~Ao,
``Potential in stochastic differential equations: novel construction,''
\textit{Journal of Physics A}, vol.~37, pp.~L25--L30, 2004.

\bibitem{RaoEsposito2016}
R.~Rao and M.~Esposito,
``Nonequilibrium thermodynamics of chemical reaction networks: wisdom
from stochastic thermodynamics,''
\textit{Physical Review X}, vol.~6, p.~041064, 2016.

\end{thebibliography}

\end{document}
